%! Mode:: "TeX:UTF-8"
%! TEX program = xelatex
\PassOptionsToPackage{quiet}{xeCJK}
\documentclass[withoutpreface,bwprint]{cumcmthesis}
\usepackage{etoolbox}
\usepackage{titlesec} %设置标题靠左
\usepackage{listings} %利用此环境插入py代码
\BeforeBeginEnvironment{tabular}{\zihao{-5}}
\usepackage[numbers,sort&compress]{natbib}  % 文献管理宏包
\usepackage[framemethod=TikZ]{mdframed}  % 框架宏包
\usepackage{url}  % 网页链接宏包
\usepackage{pdfpages} %插入附录部分
% 三线表格式
\usepackage{booktabs}
\usepackage{adjustbox}
\usepackage{subcaption}  % 子图宏包
\newcolumntype{C}{>{\centering\arraybackslash}X}
\newcolumntype{R}{>{\raggedleft\arraybackslash}X}
\newcolumntype{L}{>{\raggedright\arraybackslash}X}

% 定义部分section标题靠左的命令
\newcommand{\leftalignedsection}[1]{%
    \section*{#1}%
    \addcontentsline{toc}{section}{#1}%
}

% 使用titlesec包设置部分section标题靠左
\titleformat{name=\section,numberless}[block]{\raggedright\bfseries\Large}{}{0pt}{}

\title{生产线的故障自动识别与人员配置}  % 论文标题
\tihao{}  % 题号
\baominghao{}  % 报名号
\schoolname{}  % 学校
\membera{}  % 队员a
\memberb{}  % 队员b
\memberc{}  % 队员c
\supervisor{}  % 指导老师
\yearinput{}
\monthinput{}
\dayinput{}

%%%%%%%%%%%%%%%%%%%%%%%%%%%%%%%%%%%%%%%%%%%%%%%%%%%%%%%%%%%%%
%% 正文
\begin{document}

\maketitle

\begin{abstract}
工业生产线智能化控制技术在提升生产效率和产品质量方面具有突出赋能作用，但其表现受到故障智能报警技术的可靠性、人员配置的合理性以及生产规模的扩大等因素的影响，致使生产效率和经济成本呈现波动。本团队运用 \textbf{Python}进行问题数据的排除，采用\textbf{XGBoost}、 \textbf{LightGBM} 和\textbf{RandomForest}等数学模型对各装置故障的数据特征进行验证，运用决策树模型和\textbf{Gradient Boosting}实现生产线中各装置故障的自动即时报警。同时结合\textbf{SVM}及线性回归建立产品产量、合格率与生产线、操作人员等因素之间的数学模型。运用聚类分析及\textbf{MILT}模型，针对生产规模的扩展，制定出操作人员最佳排班方案。本团队所提出的数学模型在故障自动识别以及分析产品的产量、合格率与生产线、操作人员之间关系两方面具有高精确性，对于保障生产和降本增效具有重要的理论意义和实践价值。

\textbf{对于问题一,}本团队对生产线中各装置故障数据进行了深入分析。特别注意到数据集存在着类别不平衡的情况，即某些装置的故障样本数量较少，这可能会对模型的训练和泛化能力产生负面影响。本团队在模型构建过程中决策分类器（DecisionTreeClassifier）来解决这一问题，构建\(LightGBM\)模型，进行数据处理，并使用六种模型分别训练，分析得到\(LightGBM\)模型综合效果最优，较好地拟合本问题。

\textbf{对于问题二,}分别对附件二中的数据进行数据清洗、归类整合以及模型训练，使用问题一中针对各个故障建立的\(LightGBM\)模型，有效实现了生产线中各装置故障的自动即时报警，分析获得故障报警的日期、开始时间与持续时长。计算出每条生产线中每月的故障总次数，最长持续时间以及最短持续时间。

\textbf{对于问题三,}对问题、生产线及操作人员数据集的分析，首先进行数据清洗，对生产线一年内整体情况进行汇总、整合及迭代，针对题目要求和数据特点我们选择将其视为回归问题进行求解，选择均方误差\(MSE\)和决定系数（$R^2$）作为模型评估指标。基于所有特征使用\(XGBoost\)、\(SVM\)、\textbf{神经网络回归}、\textbf{随机森林}对训练数据进行拟合，采用特征值\(SHAP\)可视化进行特征值贡献排序，针对每一因素进行\(Kruskal Wallis\)检验，对拟合度进行深度分析并提出建议，作用于模型优化。

\textbf{对于问题四,} 首先结合问题三的分析结果，分析有目标函数为对最大化合格率与生产数量的乘积进行求和获得最大值，使用 \(Gradient Boosting\)、聚类分析及 \(MILP\)模型，验证工作天数、班次人数、工龄分布等约束条件与最终求和值的相关性，得到最优化操作人员排班方案。第四题着眼于现实中实际情况，需要联系现实分析相关性产生的现实原因及后续影响，探究此排班优化方案模型对于现实生活中工厂生产线排班问题具有显著指导意义。

\keywords{工业智能化\quad  故障自动报警\quad  \(LightGBM\)\quad 决策树模型 \quad 随机森林模型 \quad 特征工程 \quad \(Gradient Boosting\) \quad \(Kruskal Wallis\)检验 \quad 聚类分析}
\end{abstract}
%%%%%%%%%%%%%%%%%%%%%%%%%%%%%%%%%%%%%%%%%%%%%%%%%%%%%%%%%%%%% 

\begin{center}
    \tableofcontents  % 目录
\end{center}

\newpage

%%%%%%%%%%%%%%%%%%%%%%%%%%%%%%%%%%%%%%%%%%%%%%%%%%%%%%%%%%%%%  
\section{问题描述}

\subsection{问题背景}

随着新兴信息技术的大规模应用，工业生产线的智能化控制技术日益成熟。自动生产线 
可以自动完成物品传送、物料填装、产品包装和质量检测等过程，极大地提高了生产效率和 
产品质量，减少了生产成本。自动生产线融入故障智能报警技术，能避免因故障带来的生产 
中断和经济损失；同时合理的人员配置，能够减少资源浪费、提高生产效率。

%%%%%%%%%%%%%%%%%%%%%%%%%%%%%%%%%%%%%%%%%%%%%%%%%%%%%%%%%%%%% 
\subsection{本题所给出的信息和数据}

\textbf{（1）} 10 条生产线一年的运行记录数据与运行记录表中各字段的说明。（详见题目附件一）

\textbf{（2）} 不包含各装置故障信息的2 条生产线一年的运行记录数据（详见题目附件二）

\textbf{（3）} 10 条生产线一年的运行记录数据以及各生产线操作人员的信息（详见题目附件三）

\textbf{（4）} 最佳排班方案的模板文件（详见题目附件四）

%%%%%%%%%%%%%%%%%%%%%%%%%%%%%%%%%%%%%%%%%%%%%%%%%%%%%%%%%%%%% 
\subsection{待解决的问题}

\subsubsection{问题一} 根据附件 1 中的数据，分析生产线中各装置故障的数据特征，构建故障报警模 
型，实现故障的自动即时报警。 

\subsubsection{问题二} 应用问题 1 所建立的模型，对数据进行分析判断，实现生产线中各 装置故障的自动即时报警，给出故障报警的日期、开始时间与持续时长，并在论文中给出每月的故障总次数及最长与最短的持续时长。

\subsubsection{问题三} 根据数据，分析产品的产量、合格率与生产线、操作人员等因素的 
关系。 

\subsubsection{问题四} 根据实际情况，现需要扩大生产规模，将生产线每天的运行时间从 8 小时增加 到 24 小时不间断生产。针对问题 3 的 10 条生产线，结合问题 3 的分析结果，考虑生产线与 
操作人员的搭配，制定最佳的操作人员排班方案。
要求排班满足条件如下：

\textbf{（1）} 各操作人员做五休二，尽量连休 2 天； 

\textbf{（2）} 各操作人员每班连续工作 8 小时；

\textbf{（3）} 班次时间：早班（8:00-16:00）、中班（16:00-24:00）、晚班（0:00-8:00）；

\textbf{（4）} 各工龄操作人员的人数比例与问题 3 中的比例相同；

\textbf{（5）} 各操作人员的班次安排尽量均衡。

%%%%%%%%%%%%%%%%%%%%%%%%%%%%%%%%%%%%%%%%%%%%%%%%%%%%%%%%%%%%% 

\section{模型假设}

为简化问题，本文做出以下假设：

\begin{itemize}[itemindent=2em]

\item 假设1：假设生产线运行时间的延长不会影响故障率，因为故障通常由设备磨损和维护不当引起，而非时间因素造成的。

\item 假设2：假设生产线上的操作人员在不同时段表现稳定，他们的工作质量主要受个人能力和经验的影响，而非特定的日期或时间。

\item 假设3：假设生产线上各装置的故障发生是随机的，不受具体时间或日期的影响，可能是由于设备本身的缺陷或外部因素引起的。

\item 假设4：假设生产线运行时间的延长不会对产品质量产生直接影响，因为质量通常受制于生产过程中的其他因素，如材料和操作技术。

\item 假设5：假设产品的合格率与生产线运行时间和日期无关，而主要受到生产过程中的设备状态和操作人员技术水平的影响。

\item 假设6：假设无论是白天还是夜晚生产，产品的质量和生产线的故障率都保持稳定，因为生产线的运行与时间无关，而是受到设备状态和人员操作的影响。

\end{itemize}

%%%%%%%%%%%%%%%%%%%%%%%%%%%%%%%%%%%%%%%%%%%%%%%%%%%%%%%%%%%%% 

\section{模型介绍}

\subsection{随机森林}

\subsubsection{决策树}

决策树\textsuperscript{\cite{cramer1976estimation}}是一种树形结构，其中每个内部节点表示一个属性上的判断，每个分支代表一个判断结果的输出，最后每个叶节点代表一种分类情况。决策树的构建是通过训练数据，并根据基尼系数的增益进行统计而来
\par 基尼系数 \( Gini(D) \) 表示数据集 \( D \) 中样本的差异程度，基尼系数越大，表示数据集的种类越多样，即表示有多种的分类结果，数据集的当前特征越多样，即越不纯。其具体公式如下所示：
\[
Gini(D) = 1 - \sum_{k=1}^{|y|} \left( \frac{1}{|y|} \right)^2
\]

其中，\( |y| \) 表示类别的数量。

基尼增益 \( Gini(D, \alpha) \) 是根据每个特征作为判断依据时的基尼指数计算而来，具体公式如下：

\[
Gini(D, \alpha) = Gini\_index(D, \alpha)
\]

决策树的构建是不断地对内部节点进行分类，直至显示最终的分类结果为止。在实际工程中，需要控制决策树的深度，以防止过拟合的情况。

%%%%%%%%%%%%%%%%%%%%%%%%%%%%%%%%%%%%%%%%%%%%%%%%%%%%%%%%%%%%% 
\subsubsection{Bagging 算法定理}

\(Bagging\)算法基于自助采样法（\(bootstrap sampling\)）。给定包含 \( N \) 个样本的训练数据集 \( D \)，自助采样法是这样进行的：先从 \( D \) 中随机取出一个样本放入采样集 \( D_s \) 中，再把该样本放回 \( D_s \) 中（有放回的重复独立采样）。经过 \( N \) 次随机采样操作，得到包含 \( N \) 个样本的采样集 \( D_s \) 。

\(Bagging\)首先采用 \( M \) 轮自助采样法，获得 \( M \) 个包含 \( N \) 个训练样本的采样集。然后，基于这些采样集训练出一个基学习器。最后将这 \( M \) 个基学习器进行组合。

%%%%%%%%%%%%%%%%%%%%%%%%%%%%%%%%%%%%%%%%%%%%%%%%%%%%%%%%%%%%% 

\subsubsection{随机森林}

随机森林\textsuperscript{\cite{belgiu2016random}}（\(Random Forest\)，\(RF\)）是一种以决策树为基学习器的\(Bagging\)算法。在随机森林中构建决策树时，引入了随机属性选择的概念。具体来说，在选择节点的划分属性时，首先从该节点的属性集合中随机选择一个包含 \( k \) 个属性的子集，然后再从这个子集中选择一个最优属性用于划分。

随机森林的训练效率较高，因为每棵树构建时的样本都是由训练集经过有放回抽样得来的。另外，在树的构造过程中拆分每个节点时，可以从所有输入特征或大小为 \( \text{max\_features} \) 的随机子集中找到最佳拆分。

随机森林通过组合不同的树来减少方差，有时会以略微增加偏差为代价。在实践中，由于方差的减小通常是很明显的，因此在总体上产生了更好的模型。

%%%%%%%%%%%%%%%%%%%%%%%%%%%%%%%%%%%%%%%%%%%%%%%%%%%%%%%%%%%%% 


\subsection{ARIMA模型}

\subsubsection{ARIMA介绍}

\begin{itemize}[itemindent=2em]

\item \(ARIMA\) \textsuperscript{\cite{boschen1994long}}( \(Auto Regressive Integrated Moving Average\))是一种基于时间序列历史值和历史值上的预测误差来对当前做预测的模型。

\item \(ARIMA\)整合了自回归项\(AR\)和滑动平均项\(MA\)。

\item \(ARIMA\)可以建模任何存在一定规律的非季节性时间序列。

\item 如果时间序列具有季节性，则需要使用 \(SARIMA(Seasonal ARIMA)\)建模，后续会介绍。

\end{itemize}

%%%%%%%%%%%%%%%%%%%%%%%%%%%%%%%%%%%%%%%%%%%%%%%%%%%%%%%%%%%%% 

\subsubsection{ARIMA模型参数}

\(ARIMA\)模型有三个超参数：\(p\),\(d\),\(q\)

\textbf{（1）} \(p\):

\hspace{2em} \(AR\)(自回归)项的阶数。需要事先设定好，表示\(y\)的当前值和前\(p\)个历史值有关。

\textbf{（2）} \(d\):

\hspace{2em} 使序列平稳的最小差分阶数，一般是1阶。非平稳序列可以通过差分来得到平稳序列，但是过度的差分，会导致时间序列失去自相关性，从而失去使用\(AR\)项的条件。

\textbf{（3）} \(q\):

\hspace{2em} \(MA\)(滑动平均)项的阶数。需要事先设定好，表示\(y\)的当前值和前\(q\)个历史值\(AR\)预测误差有关。实际是用历史值上的\(AR\)项预测误差来建立一个类似归回的模型。

%%%%%%%%%%%%%%%%%%%%%%%%%%%%%%%%%%%%%%%%%%%%%%%%%%%%%%%%%%%%% 

\subsubsection{ARIMA模型表示}

\textbf{（1）} \(AR\)项表示:

\hspace{2em} 一个\(p\)阶的自回归模型可以表示如下：
\[y_{t} = c + \phi _{1}y_{t-1} +  \phi _{2}y_{t-2} + ...+ \phi _{p}y_{t-p}+\varepsilon _{t} \]

c是常数项，$\varepsilon _{t}$是随机误差项。

对于一个\(AR\)模型而言：

\begin{itemize}[itemindent=2em]

\item 当$\phi _{1}$=0时,$y_{t}$相当于白噪声；

\item 当$\phi _{1}$=1并且\(c\)=0时,$y_{t}$相当于随机游走模型；

\item 当$\phi _{1}$=1并且c≠0时,$y_{t}$相当于带漂移的随机游走模型；

\item 当$\phi _{1}$<0时,$y_{t}$倾向于在正负值之间上下浮动。

\item 如果时间序列具有季节性，则需要使用\(SARIMA(Seasonal ARIMA)\)建模，后续会介绍。

\end{itemize}

\textbf{（2）} \(MA\)项表示:

\hspace{2em} 一个\(q\)阶的预测误差回归模型可以表示如下：
\[y_{t} = c + \varepsilon _{t}+\theta_{1}  \varepsilon _{t-1}+ \theta_{2}  \varepsilon _{t-2} +...+\theta_{q}  \varepsilon _{t-q}\]

\(c\)是常数项，$\varepsilon _{t}$是随机误差项。

$y_{t}$是历史预测误差的加权移动平均值，\(q\)指定历史预测误差的期数。

\textbf{（3）} 完整表示:
\[Y_{t} = c + \phi _{1}Y_{t-1} +  \phi _{2}Y_{t-2} + ...+ \phi _{p}Y_{t-p}+\varepsilon _{t}+\varepsilon _{t}+\theta_{1}  \varepsilon _{t-1}+ \theta_{2}  \varepsilon _{t-2} +...+\theta_{q}  \varepsilon _{t-q}\]

即: 被预测变量\(Yt\) = 常数+\(Y\)的\(p\)阶滞后的线性组合 + 预测误差的\(q\)阶滞后的线性组合

%%%%%%%%%%%%%%%%%%%%%%%%%%%%%%%%%%%%%%%%%%%%%%%%%%%%%%%%%%%%% 
\subsubsection{ARIMA模型定阶}

差分阶数\(d\):

\begin{itemize}[itemindent=2em]

\item 如果时间序列本身就是平稳的，就不需要差分，所以此时\(d\)=0。

\item 如果时间序列不平稳，就需要看时间序列的 \(acf\) 图，如果\(acf\) 表现为 10 阶或以上的拖尾，那么需要进一步的差分，如果\(acf\) 表现为 1 阶截尾，则可能是过度差分，最好的差分阶数是使\(acf\) 先拖尾几阶，然后截尾。

\item 有些情况下，可能在 2 个阶数之间无法确定用哪个，因为\(acf\)的表现近似，此时就选择标准差小的序列。

\item 下面是原时间序列、一阶差分后、二阶差分后的\(acf\)图：

\begin{figure}[htbp]
  \centering
  \includegraphics[width=0.75\textwidth]{image_1.png}
  \caption{acf图分析}
  \label{fig:example_1}  %用于文中的图片引用
\end{figure}

\end{itemize}

可以看到，原序列的\(acf\)图的拖尾阶数过高，而二阶差分后的截尾阶数过小，所以一阶差分更合适。

%%%%%%%%%%%%%%%%%%%%%%%%%%%%%%%%%%%%%%%%%%%%%%%%%%%%%%%%%%%%% 
\subsection{XGBOOST}

\(XGBoost\) \textsuperscript{\cite{sharma2024xgboost}}是一种实现\(Boosting Tree\)的算法，它利用二阶泰勒展开对损失函数进行近似，并结合一阶和二阶导数信息。此外，\(XGBoost\)还引入了适用于树模型的正则化项，用于控制模型的复杂度。其正则化项包括树的叶子节点个数以及每个叶子节点输出分数的 \(L2\) 平方和。

目标函数为：
\[
\tilde{L}^{(t)} = \sum^n_{i=1} \left[ g_i w_q(x_i) + \frac{1}{2} h_i w_q^2(x_i) \right] + \gamma T + \frac{1}{2} \lambda \sum^T_{j=1} \left[ \left( \sum_{i\in I_j} g_i \right) w_j + \frac{1}{2} \left( \sum_{i\in I_j} h_i + \lambda \right) w_j^2 \right] + \gamma T
\]

其中，\( \tilde{L}^{(t)} \) 表示第 \( t \) 轮迭代后的目标函数，\( g_i \) 和 \( h_i \) 分别表示样本 \( i \) 的一阶和二阶导数，\( w \) 表示叶子节点的输出分数，\( \gamma \) 和 \( \lambda \) 是正则化参数，\( T \) 是叶子节点的个数。

通过对目标函数求偏导，并令偏导数等于零，可以求解得到最优的叶子节点输出分数 \( \omega^* \)：

\[
\omega^* = -\frac{\sum_{i\in I_j} g_i}{\sum_{i\in I_j} h_i + \lambda}
\]

在选取分裂点时，也以最小化目标函数为目标。假设在某次选取分裂点进行划分后，左右子节点分别为 \( I_L \) 和 \( I_R \)，且 \( I = I_L \cup I_R \)，那么损失函数减小的值为：

\[
L_{split} = \frac{1}{2} \left[ \frac{\left( \sum_{i\in I_L} g_i \right)^2}{\sum_{i\in I_L} h_i + \lambda} + \frac{\left( \sum_{i\in I_R} g_i \right)^2}{\sum_{i\in I_R} h_i + \lambda} + \frac{\left( \sum_{i\in I_j} g_i \right)^2}{\sum_{i\in I_j} h_i + \lambda} \right] - \gamma
\]

其中，\( L_{split} \) 表示在分裂点处的损失函数减小值，\( \gamma \) 是正则化参数。

%%%%%%%%%%%%%%%%%%%%%%%%%%%%%%%%%%%%%%%%%%%%%%%%%%%%%%%%%%%%% 

\subsection{关联规则数据挖掘}

\subsubsection{问题定义}

\begin{itemize}[itemindent=2em]

\item 定义 1--\textbf{项集}：设I={$i_{1}$，$i_{2}$，…，$i_{n}$}是\(n\)个不同的数据项组成的集合。包含 0 个或多个数据项的集合，称之为项集。

\item 定义 2-- \textbf{k-项集}：如果一个项集包含\(k\)个数据项，则称之为\(k\)-项集。

\item 定义 3-- \textbf{事务}：事务是不同的数据项组成的集合。每一条事务\(T\)是\(I\)的一个子集，即\(T\) $\subseteq$ \(I\)。

\item 定义 4-- \textbf{事务数据集}：事务数据集\(D\)={$T_{1}$，$T_{2}$，…，$T_{m}$}是\(m\)条事务组成的集合。

\item 定义 5-- \textbf{项集的支持度计数}：项集\(X\)在\(D\)中的支持度计数表示\(D\)中包含\(X\)的事务数，其形式化定义如下:
\[count(X)=|{ T_{i} |X\subseteq  T_{i} , T_{i}\in D}|\]

其中，||表示集合中元素的个数。

\item 定义6-- \textbf{项集的支持度}：项集\(X\)在\(D\)中的支持度表示包含\(X\)的事务在\(D\)中所占的比例，其形式化定义如下:
\[support(X)= \frac{|{T_{i}|X\subseteq T_{i} , T_{i}\in D}|} {|D|} \]

\item 定义7-- \textbf{频繁项集}：如果项集\(X\)的支持度不小于用户给定的最小支持度阈值 minSup，则称\(X\)为频繁项集。

\item 定义8-- \textbf{候选项集}：可能成为频繁项集的项集称为候选项集。

\end{itemize}

%%%%%%%%%%%%%%%%%%%%%%%%%%%%%%%%%%%%%%%%%%%%%%%%%%%%%%%%%%%%% 

\subsubsection{问题分析}

如果一个事务数据集中包含 5 个{\(a\)}，{\(b\)}，{\(c\)}，{\(d\)}，{\(e\)}不同的数据项，则该问题的解空间包含 25 个候选项集，如下图所示。

\begin{figure}[htbp]
  \centering
  \includegraphics[width=0.75\textwidth]{image_2.png}
  \caption{问题的解空间分析}
  \label{fig:example_2}  %用于文中的图片引用
\end{figure}

为了计算该解空间中每一个候选项集的支持度，\(Brute-force\)方法通过扫描事务数据集，将所有候选项集与事务逐一比较。如果事务中包含该候选项集，则增加候选项集的支持度。该过程如下图 \ref{fig:example_3} 所示。

\begin{figure}[h]
  \centering
  \includegraphics[width=0.75\textwidth]{image_3.png}
  \caption{计算解空间过程分析}
  \label{fig:example_3}  %用于文中的图片引用
\end{figure}

从而，可知\(Brute-force\)方法的时间复杂度约为\(O(NMw)\)，其中\(N\)是事务数据集的大小，\(M\)是候选项集的数量，\(w\)是平均事务的长度。由于\(M=2k\)，\(k\)是数据项的数量，因此，该方法的时间复杂度是指数级别。

%%%%%%%%%%%%%%%%%%%%%%%%%%%%%%%%%%%%%%%%%%%%%%%%%%%%%%%%%%%%% 

\newpage
\subsubsection{Apriori算法分析}

\textbf{（1）} 基本思想:

\(Apriori\) \textsuperscript{\cite{bhargava2013association}}算法采用了逐层搜索的策略来对解空间进行遍历。在遍历的过程中 ，该算法采用了先验原理（如果一个项集是频繁项集，则其任意子集均是频繁项集。）来对解空间进行剪枝，减少候选项集的数量。

\begin{figure}[htbp]
  \centering
  \includegraphics[width=0.75\textwidth]{image_4.png}
  \caption{解的空间剪枝}
  \label{fig:example_4}  %用于文中的图片引用
\end{figure}

给定如上图所示的解空间，如果候选项集 {A，B}不是频繁项集，则该候选项集的任意超集均不可能成为频繁项集，因此，无需计算这些项集的支持度，可以将其从解空间中剪枝，减少不必要的计算量。

%%%%%%%%%%%%%%%%%%%%%%%%%%%%%%%%%%%%%%%%%%%%%%%%%%%%%%%%%%%%% 

\textbf{（2）}算法过程

\(Apriori\)算法首先扫描一次事务数据集，计算各个数据项的支持度，从而得到频繁 1-项集。将这些频繁 1-项集自连接来生成候选 2-项集。通过再一次扫描事务数据集，计算各个候选项集的支持度，从而得到频繁 2-项集。\(Apriori\)算法按照这个方式不断迭代，直至不再产生新的候选项集或频繁项集为止。其中，将两个 k-项集自连接来生成候选 (k+1)-项集的要求是这两个 k-项集除了最后一个数据项不同，其余数据项均相同。所生成的不满足先验原理的候选 (k+1)-项集将被删除。

\begin{figure}[htbp]
  \centering
  \includegraphics[width=0.75\textwidth]{image_5.png}
  \caption{\(Apriori\)算法过程}
  \label{fig:example_5}  %用于文中的图片引用
\end{figure}

%%%%%%%%%%%%%%%%%%%%%%%%%%%%%%%%%%%%%%%%%%%%%%%%%%%%%%%%%%%%% 

\subsection{线性回归模型}

\subsubsection{线性回归定义}

线性回归模型是一种用于描述因变量 \( y \) 与一个或多个自变量 \( X = (X_1, X_2, ..., X_p)^T \) 之间线性关系的统计模型。其数学表达式如下：
\[ y = \beta_0 + \beta_1 X_1 + \beta_2 X_2 + ... + \beta_p X_p + \epsilon \]

其中，\( \beta_0 \) 是截距（或偏置），\( \beta_1, \beta_2, ..., \beta_p \) 是特征 \( X_1, X_2, ..., X_p \) 的系数，\( \epsilon \) 是误差项。

\subsubsection{线性回归基本原理}

线性回归模型的基本原理是通过拟合一个线性方程来描述自变量 \( X \) 和因变量 \( y \) 之间的关系，使得该方程能够最好地拟合已知的观测数据。其核心步骤包括：

\textbf{（1）} 建立线性关系：线性回归假设因变量 \( y \) 与自变量 \( X \) 之间存在线性关系，即可以被表示为 \( X \) 的线性组合。

\textbf{（2）} 最小化残差平方和：通过调整模型参数 \( \beta_0, \beta_1, ..., \beta_p \) 来使模型预测值与实际观测值之间的差异最小化。这通常通过最小化残差平方和来实现，即最小化残差 \( \epsilon \) 的平方和。

\textbf{（3）} 求解参数：通过最小化残差平方和，可以使用不同的优化算法（如梯度下降、最小二乘法等）来求解模型参数 \( \beta_0, \beta_1, ..., \beta_p \)。

\textbf{（4）} 模型评估：完成参数估计后，需要对模型进行评估以确保模型的准确性和可靠性。常用的评估指标包括均方误差（\(Mean Squared Error\), \(MSE\)）、\(R\)平方（\(Coefficient\) \(of\) \(Determination\), \( R^2 \)）等。

%%%%%%%%%%%%%%%%%%%%%%%%%%%%%%%%%%%%%%%%%%%%%%%%%%%%%%%%%%%%% 

\subsection{梯度提升树}

梯度提升树模型是一种基于回归树的集成学习方法，它通过构造多个弱的回归树作为基学习器，并将这些树的结果累加起来作为最终的预测输出。由于其简单的特征处理方式和出色的效果，梯度提升树模型被认为是传统统计学习中效果最好的方法之一。

梯度提升树算法是集成学习\(Boosting\)家族的一员，它在训练时采用前向分步算法。首先，确定第一棵树拟合的值，然后基于之前所有树的误差来更新并训练下一棵树，一步一步迭代下去直到梯度提升树模型构建完毕。在训练过程中，首先确定初始提升树 \( f(x) = 0 \)，然后在后续训练时，第 \( m \) 步的模型表示为：
\[ f_m(x) = f_{m-1}(x) + T(x;\Theta_m) \]

其中，\( f_{m-1}(x) \) 是当前模型，通过经验风险最小化确定下一棵树的参数 \( \Theta_m \)。

通过梯度下降法对损失函数进行求解，可以得到梯度提升树的算法如下：

\textbf{（1）} 输入：训练数据集 \( T = \{(x_1, y_1), (x_2, y_2), ..., (x_N, y_N)\} \), \( x_1 \in X \subseteq R^n \), \( y_1 \in Y \subseteq R^n \), 损失函数 \( \text{Cost}(y, f(x)) \)。

\textbf{（2）} 过程：

\begin{itemize}[itemindent=2em]

\item 初始化：\( f_0(x) = \text{arg}\min \sum^N_{i=1} L(y, c) \)；

\item 对 \( m = 1, 2, ..., M \)：

(a)对 \( i = 1, 2, ..., N \)，计算：\( \gamma_{mi} = -\left[\frac{\partial L(y_i, f(x_i))}{\partial f(x_i)}\right]f(x) = f_{m-1}(x) \)；

(b)对 \( \gamma_{mi} \) 拟合一个回归树，得到第 \( m \) 棵树的叶节点区域 \( R_{mj} \)，\( j = 1, 2, ..., J \)；

(c)对 \( j = 1, 2, ..., J \)，计算：\( c_{mj} = \text{arg}\min \sum_{x_i \in R_{mj}} L(y_i, f_{m-1}(x_i) + c) \)；

(d)更新 \( f_m(x) = f_{m-1}(x) + \sum^J_{j=1} c_{mj} I(x \in R_{mj}) \)；

\item 得到回归树：\( \hat{f}(x) = f_M(x) = \sum^M_{m=1} \sum^J_{j=1} c_{mj} I(x \in R_{mj}) \)。

\end{itemize}

\textbf{（3）} 输出：回归树 \( \hat{f}(x) \)。

%%%%%%%%%%%%%%%%%%%%%%%%%%%%%%%%%%%%%%%%%%%%%%%%%%%%%%%%%%%%% 

\subsection{LASSO模型}

\(LASSO\) \textsuperscript{\cite{rasmussen2012tutorial}}（最小绝对收缩与选择算子）是一种压缩估计方法，最早由\(Robert\) Tibshirani于 1996 年提出。该方法通过构造惩罚函数来获得一个更加简洁的模型，从而压缩一些回归系数。其核心思想是限制回归系数的绝对值之和小于某个固定值，并将一些回归系数设定为零。这样做的优势在于保留了子集收缩的特点，适用于处理具有复共线性数据的有偏估计问题。

\(LASSO\)模型是在最小化残差平方和（\(RSS\)）的基础上加入了 \( l_p \) 范数作为惩罚约束。 \( l_p \) 范数的优势在于，当惩罚系数 \( \lambda \) 充分大时，可以将某些待估系数精确地压缩到零，从而实现变量选择的功能。

通过交叉验证法，我们可以选择合适的 \( \lambda \) 值。具体做法是，对于给定的 \( \lambda \) 值，进行交叉验证，选取交叉验证误差最小的 \( \lambda \) 值，然后利用全部数据重新拟合模型。

这样的过程可以帮助我们在模型中找到合适的正则化参数，从而得到更稀疏的模型，进而提高模型的泛化能力和解释性。

%%%%%%%%%%%%%%%%%%%%%%%%%%%%%%%%%%%%%%%%%%%%%%%%%%%%%%%%%%%%% 

\subsection{LightGBM模型}

\(LightGBM\) \textsuperscript{\cite{nam2023apply}}是一个梯度\(Boosting\)框架，使用基于决策树的学习算法。它可以说是分布式的，高效的，有以下优势：

\textbf{（1）} 更快的训练效率；

\textbf{（2）} 低内存使用；

\textbf{（3）} 更快的准确率；

\textbf{（4）} 支持并行化学习。

%%%%%%%%%%%%%%%%%%%%%%%%%%%%%%%%%%%%%%%%%%%%%%%%%%%%%%%%%%%%% 

\subsubsection{GOSS介绍}

\(GOSS\)（Gradient-based One-Side Sampling）通过区分不同梯度的实例，保留较大梯度实例同时对较小梯度进行随机采样的方式，从而减少计算量，提高效率。在提升树训练过程中，目标函数学习的是负梯度（近似残差）。

\(GOSS\)的计算步骤如下：首先，选取前 \( a \% \) 个较大梯度的值作为大梯度值的训练样本；接着，从剩余的 \( 1 - a \% \) 个较小梯度的值中，随机选取其中的 \( b \% \) 个作为小梯度值的训练样本。对于较小梯度的样本，即 \( b \% \times (1 - \frac{1}{100}) \times samples \) 个样本，计算信息增益时将其放大 \( \frac{(1 - a)}{b} \) 倍。总的来说，用 \( a \% \times samples + b \% \times (1 - a \%) \times samples \) 个样本作为训练样本。这样的构造旨在尽可能保持与总的数据分布一致，并确保小梯度值的样本得到充分训练。

%%%%%%%%%%%%%%%%%%%%%%%%%%%%%%%%%%%%%%%%%%%%%%%%%%%%%%%%%%%%% 

\subsubsection{独立特征合并}

独立特征合并（\(EFB\)）通过特征捆绑的方式减少特征维度，提高计算效率。通常情况下，被捆绑的特征都是互斥的，即一个特征值为零，另一个特征值不为零。如果两个特征并不完全互斥，可以用一个指标衡量特征不互斥程度，称之为冲突比率。当冲突比率较小时，可以选择将不完全互斥的特征捆绑，而不影响最后的精度。

%%%%%%%%%%%%%%%%%%%%%%%%%%%%%%%%%%%%%%%%%%%%%%%%%%%%%%%%%%%%% 

\subsubsection{特征融合}

特征融合构建一个以特征为图中的点 \( V \)，以特征之间的总冲突为图中的边 \( E \)。冲突值是特征之间的 \(cos\) 夹角值，以尽可能保证特征之间非零元素不在同一个行上。寻找合并特征且使得合并的簇个数最小，这是一个图着色问题。因此，需要使用近似的贪心算法来完成。首先按照度对每个点（特征）进行降序排序，然后将特征合并到冲突数小于 \( K \) 的簇或者新建另外一个簇。

%%%%%%%%%%%%%%%%%%%%%%%%%%%%%%%%%%%%%%%%%%%%%%%%%%%%%%%%%%%%% 

\subsubsection{特征合并}

特征合并将这些簇中的特征合并起来，并重新构建合并后簇特征的范围。在第一个\(for\)循环中，记录每个特征与之前特征累积总范围。在第二个\(for\)循环中，根据之前的\(bin Ranges\)重新计算出新的\(bin value\)，保证特征之间的值不会冲突。这是针对稀疏矩阵进行的优化。由于之前\(Greedy Bundling\)算法对特征进行冲突检查，保证簇内特征冲突尽可能少，因此特征之间的非零元素不会有太多的冲突。

EFB的算法步骤如下：

\textbf{（1）} 将特征按照非零值的个数进行排序；

\textbf{（2）} 计算不同特征之间的冲突比率；

\textbf{（3）} 遍历每个特征并尝试合并特征，使冲突比率最小化。我们称使用\(GOSS\)算法和\(EFB\) 算法的梯度提升树（\(GBDT\)）为 \(LightGBM\)。

%%%%%%%%%%%%%%%%%%%%%%%%%%%%%%%%%%%%%%%%%%%%%%%%%%%%%%%%%%%%% 

\subsection{Adaboost}

\(Adaboost\)是一种迭代算法，其核心思想是针对同一个训练集训练不同的分类器(弱分类器)，然后将这些弱分类器集合起来，构成一个更强的最终分类器（强分类器）。后一个模型的训练永远是在前一个模型的基础上完成。

这里的集合起来的策略是通过提高前一轮分类器分类错误的样本的权值，降低分类分类正确的样本权值，对于那些没有被分类正确的样本会得到后面分类器更多的关注。然后可以产生很多的弱分类器，通过多数加权投票组合这些弱分类器，加大误差率小的分类器，减少误差率大的分类器，使其在表决中起到较少的作用。

\subsubsection{算法思想}

\textbf{（1）} 初始化训练样本的权值分布，每个样本具有相同权重；

\textbf{（2）} 训练弱分类器，如果样本分类正确，则在构造下一个训练集中，它的权值就会被降低；反之提高。用更新过的样本集去训练下一个分类器；

\textbf{（3）} 将所有弱分类器组合成强分类器，各个弱分类器的训练过程结束后，加大分类误差率小的弱分类器的权重，降低分类误差率大的弱分类器的权重。

我们的步骤可以如下图：

\begin{figure}[htbp]
  \centering
  \includegraphics[width=0.75\textwidth]{image_6.png}
  \caption{\(Adaboost\)步骤分析}
  \label{fig:example_6}  %用于文中的图片引用
\end{figure}

由\(Adaboost\)算法的描述过程可知，该算法在实现过程中根据训练集的大小初始化样本权值，使其满足均匀分布，在后续操作中通过公式来改变和规范化算法迭代后样本的权值。样本被错误分类导致权值增大，反之权值相应减小，这表示被错分的训练样本集包括一个更高的权重。这就会使在下一轮时训练样本集更注重于难以识别的样本，针对被错分样本的进一步学习来得到下一个弱分类器，直到样本被正确分类。在达到规定的迭代次数或者预期的误差率时，则强分类器构建完成。

\subsubsection{Adaboost 分类算法原理}

假设一个二分类训练样本集：
\[ T = \{(x_1, y_1), (x_2, y_2), ..., (x_m, y_m)\} \]

训练集的在第 \( k \) 个弱学习器的输出权重为：
\[ D(k) = (w_{k1}, w_{k2}, ..., w_{km});\quad w_{1i} = \frac{1}{m};\quad i = 1, 2, ..., m \]

第 \( k \) 个弱分类器 \( G_k(x) \) 在训练集上的加权误差率为：
\[ e_k = P(G_k(x_i) \neq y_i) = \sum_{i=1}^m w_{ki} I(G_k(x_i) \neq y_i) \]

第 \( k \) 个弱分类器 \( G_k(x) \) 的权重系数为：
\[ \alpha_k = \frac{1}{2} \log \frac{1 - e_k}{e_k} \]


从上式可以看出，如果分类误差率 \( e_k \) 越大，则对应的弱分类器权重系数 \( \alpha_k\) 越小，即分类误差率越小的弱分类器权重越大，这是因为对于分类误差率小的分类器，其分类能力较强，我们希望其权重较大，反之，分类误差率大的分类器，我们希望其权重较小。

每一轮迭代过程中，\(Adaboost\)算法都会更新训练样本的权值分布，使得上一轮分类错误的样本的权值增大，分类正确的样本的权值减小。其更新规则为：
\[ w_{k+1,i} = \frac{w_{ki} \cdot \exp(-\alpha_k y_i G_k(x_i))}{\sum_{i=1}^m w_{ki}} \]

这个过程是一个指数衰减的过程，对于误分的样本，其权值会增大，而对于分类正确的样本，其权值会减小。

通过上述迭代过程，我们可以得到一系列的弱分类器 \( G_k(x) \) 和对应的权重系数 \( \alpha_k \)，然后将这些弱分类器按照其权重系数进行线性加权组合，得到最终的强分类器：
\[ f(x) = \sum_{k=1}^K \alpha_k G_k(x) \]

%%%%%%%%%%%%%%%%%%%%%%%%%%%%%%%%%%%%%%%%%%%%%%%%%%%%%%%%%%%%% 

\subsection{GBDT}

梯度提升树（\(Gradient Boosting Decision Tree\)，简称\(GBDT\)）作为一种集成学习算法，在数据挖掘和机器学习领域中具有广泛的应用。本团队将深入探讨\(GBDT\)的基本原理，包括其定义、工作机制以及如何通过迭代减少误差来提高模型性能。

\subsubsection{GBDT 的定义}

\(GBDT\)的核心思想是将多个弱学习器（通常是决策树）组合成一个强大的预测模型。具体而言，\(GBDT\)的定义如下：

\begin{itemize}[itemindent=2em]

\item 初始化：首先，\(GBDT\)使用一个常数（通常是目标变量的平均值）作为初始预测值。这个初始预测值代表了我们对目标变量的初始猜测。

\item 迭代训练：\(GBDT\)是一个迭代算法，通过多轮迭代来逐步改进模型。在每一轮迭代中，\(GBDT\)都会训练一棵新的决策树，目标是减少前一轮模型的残差（或误差）。残差是实际观测值与当前模型预测值之间的差异，新的树将学习如何纠正这些残差。

\item 集成：最终，\(GBDT\)将所有决策树的预测结果相加，得到最终的集成预测结果。这个过程使得模型能够捕捉数据中的复杂关系，从而提高了预测精度。

\end{itemize}

\(GBDT\)的核心原理在于不断迭代，每一轮迭代都尝试修正前一轮模型的错误，逐渐提高模型的预测性能。

\subsubsection{GBDT的工作机制}

\textbf{(1)} 初始化：在训练开始时，\(GBDT\)使用一个初始预测值来代表整体数据的平均情况。这个初始预测值可以是目标变量的均值，也可以是其他合适的初始值。初始预测值代表了模型对整体数据的初始估计。

\textbf{(2)} 迭代训练：\(GBDT\)是一个迭代算法，通常包括多轮迭代。在每一轮迭代中，模型都会训练一棵新的决策树，这棵树的目标是减少前一轮模型的残差。具体步骤如下：

\begin{itemize}

    \item 计算残差：在每轮迭代开始时，计算当前模型对训练数据的预测值与实际观测值之间的残差。这个残差代表了前一轮模型未能正确预测的部分。

    \item 训练新的决策树：使用计算得到的残差作为新的目标变量，训练一棵新的决策树。这棵树将尝试纠正前一轮模型的错误，以减少残差。

    \item 更新模型：将新训练的决策树与之前的模型进行组合。具体地，将新树的预测结果与之前模型的预测结果相加，得到更新后的模型。
    
\end{itemize}

\textbf{(3)} 集成预测：在进行了多轮迭代训练后，\(GBDT\)将所有决策树的预测结果进行累加，得到最终的集成预测结果。这个集成过程充分利用了每棵树的贡献，使得模型能够更好地拟合数据并提高预测精度。\(GBDT\)的工作机制保证了模型的预测能力逐渐提高，每一轮迭代都尝试修正前一轮模型的错误。这种集成学习的方法使得\(GBDT\)在各种数据挖掘任务中表现出色。

\subsubsection{GBDT的优势}

\textbf{(1)} 高精度预测能力

\begin{itemize}
    \item 集成学习的威力：\(GBDT\)采用集成学习的方法，将多个弱学习器（通常是决策树）组合成一个强大的模型。这种集成学习的方式能够有效地减少模型的偏差和方差，从而提高了预测的准确性。

    \item 处理非线性关系：\(GBDT\)的非线性建模能力使其能够更好地拟合数据，提高了模型的预测准确性。
    
\end{itemize}

\textbf{(2)} 对各种类型数据的适应性

\begin{itemize}
    \item 处理混合数据类型：\(GBDT\)能够直接处理混合数据，无需将其转换成统一的格式，简化了数据预处理的步骤。

    \item 不需要特征缩放：\(GBDT\)不需要对特征进行缩放或归一化，简化了数据预处理的负担。
    
\end{itemize}

\textbf{(3)} 在数据不平衡情况下的优势

\begin{itemize}
    \item 加权损失函数：\(GBDT\)使用的损失函数允许对不同类别的样本赋予不同的权重，提高了对不平衡数据的处理能力。

    \item 逐步纠正错误：\(GBDT\)的迭代训练方式使其能够逐步纠正前一轮模型的错误，在处理不平衡数据时表现出色。
    
\end{itemize}

%%%%%%%%%%%%%%%%%%%%%%%%%%%%%%%%%%%%%%%%%%%%%%%%%%%%%%%%%%%%% 

\newpage
\section{相关理论介绍}

\subsection{缺失值处理}

\subsubsection{缺失值来源}

数据的缺失是难以避免的，可能出现的原因有多种。这些原因主要分为以下三类：

\textbf{(1)} 无意的：信息可能因为疏忽或遗漏而缺失，比如由于工作人员的疏忽或数据采集器故障等；

\textbf{(2)} 有意的：某些数据集规定将缺失值视为一种特殊的特征值；

\textbf{(3)} 不存在的：有些特征属性本身就是不存在的，比如未婚者的配偶姓名或孩子的收入状况。

\subsubsection{数据缺失类型}

在处理缺失数据之前，了解数据缺失的机制和形式至关重要。根据缺失的分布，缺失可以分为以下三种类型：

\textbf{(1)} 完全随机缺失（\(MCAR\)）：数据的缺失是完全随机的，与任何变量无关，不影响样本的无偏性；

\textbf{(2)} 随机缺失（\(MAR\)）：数据的缺失与其他变量有关，但不依赖于缺失变量本身；

\textbf{(3)} 非随机缺失（\(MNAR\)）：数据的缺失与缺失变量本身的取值有关。

针对不同类型的缺失，需要采取不同的处理方法。对于随机缺失，可以通过已知变量对缺失值进行估计，而非随机缺失的处理则相对困难。

\subsubsection{缺失值处理}

以下是处理缺失值的三种方法：删除记录，数据填补，和不处理。

\textbf{(1)} 删除记录：

\begin{itemize}
    \item 优点：这种方式简单直接。
    \item 缺点：删除记录会导致丢失大量数据，可能包含了重要信息。尤其是当缺失数据比例较大，或者缺失数据不是随机分布时，直接删除可能会使数据的分布发生偏移，例如，原本的正态分布可能变为非正态分布。这种方法通常在样本数据量很大且缺失值很少的情况下才适用，而对于样本量不大或缺失值较多的情况则不建议使用。
\end{itemize}

\textbf{(2)} 数据填补：

对缺失值的填补方法大致可以分为两种：替换和拟合，同时还有一种是使用虚拟变量。替换是通过分析数据中非缺失数据的相似性来进行填补，其核心思想是找出相似群体的共同特征；而拟合则是利用其他特征建模来填补缺失值。以下是一些常见的填补方法：

\begin{itemize}
    \item 均值填补：对于定性数据，可以使用众数填补；对于定量（比例）数据，可以使用平均值或中位数填补。通常情况下，当特征分布为正态分布时，使用平均值效果较好；而在存在异常值而不是正态分布的情况下，使用中位数效果较好。尽管这种方法简单易行，但不够准确，可能会引入噪声或改变特征原有的分布。
    
    \item 热卡填补（\(Hot deck imputation\)）：该方法通过在完整数据中寻找与缺失数据最相似的对象来进行填补。然而，不同的问题可能会采用不同的相似性标准，而且难以定义相似性标准，因此具有较多的主观因素。

    \item K最近邻法（\(K-nearest neighbors\)）：该方法利用无监督机器学习的聚类方法，通过对所有样本进行聚类划分，然后利用每个类的均值来填补各自类中的缺失值。
\end{itemize}

\textbf{(3)} 不处理：

不处理缺失值可能会导致新的噪声引入数据中，从而产生错误的结果。因此，在许多情况下，希望在保持原始信息不变的前提下处理信息系统。在实际应用中，一些模型可以处理具有缺失值的数据，比如\(XGBoost\)、\(RF\)等高级模型。

%%%%%%%%%%%%%%%%%%%%%%%%%%%%%%%%%%%%%%%%%%%%%%%%%%%%%%%%%%%%% 

\subsection{数据不平衡}

数据不平衡，也称为数据倾斜，在实际应用中是一个常见的问题，特别是在分类问题中，不同标签的样本比例往往不均衡。因此，为了消除这种不平衡，通常采用以下三种方法对数据进行采样：

\subsubsection{过采样}
过采样（\(RandomOverSampler\)）是通过主动获取更多比例较少的样本数据来解决不平衡问题。可以采用重复、自举或合成少数类样本等方法（如\(SMOTE\)）来生成新的稀有样本。然而，直接简单复制重复样本可能会导致过拟合问题，因此改进的过采样方法通常会在少数类中加入随机噪声或干扰数据，或者通过一定规则产生新的合成样本（即数据增强）。

\subsubsection{欠采样}
欠采样（\(RandomUnderSampler\)）是在数据量足够的情况下，通过减少比例较大的样本数据来平衡数据集，同时保留比例较小的样本数据。然而，欠采样的缺点是可能会丢失多数类中的一些重要信息。

\subsubsection{SMOTE采样}
\(SMOTE\)\textsuperscript{\cite{fernandez2018smote}}（\(Synthetic Minority Over-sampling Technique\)）采样的原理是在少数类样本之间进行插值来产生额外的样本。具体来说，对于一个少数类样本，首先使用近邻法找出离它最近的k个少数类样本，然后从这\(k\)个近邻点中随机选取一个，使用下列公式生成新样本：

\[ \mathbf{x}_{\text{new}} = \mathbf{x}_i + (\hat{\mathbf{x}}_i - \mathbf{x}_i) \times \delta \]

其中，\( \hat{\mathbf{x}}_i \) 是选出的k近邻点，\( \delta \in [0, 1] \) 是一个随机数。\(SMOTE\)会随机选取少数类样本用以合成新样本，而不考虑周边样本的情况。然而，这可能会导致两个问题：

\textbf{（1）} 如果选取的少数类样本周围也都是少数类样本，则新合成的样本不会提供太多有用信息；

\textbf{（2）} 如果选取的少数类样本周围都是多数类样本，这类的样本可能是噪音，则新合成的样本会与周围的多数类样本产生大部分重叠，致使分类困难。


%%%%%%%%%%%%%%%%%%%%%%%%%%%%%%%%%%%%%%%%%%%%%%%%%%%%%%%%%%%%% 

\subsection{特征提取}


特征选择是从原始数据中选取一个子集，与特征提取的区别在于不是构造新的特征，而是从现有特征中选取，构建模型而选择相关特征子集的过程。其前提是训练数据中包含大量冗余或无关特征，移除这些特征不会丢失信息。通常应用于特征较多但样本相对较少的情况。

特征选择包括产生过程、评价函数、停止准则和验证过程。首先，需要生成特征或特征子集候选集合；其次，需要衡量特征或特征子集的重要性或好坏程度，通常量化特征变量和目标变量之间的联系以及特征之间的相互联系；为了避免过拟合，常采用交叉验证评估特征的好坏；为减少计算复杂度，可能设定一个评价函数值阈值，达到阈值后停止搜索；最后，在验证数据集上验证选出的特征子集的有效性。

特征选择的方法主要分为三大类：过滤式方法、包裹式方法和嵌入式方法。

\textbf{（1）} 过滤式方法：利用统计指标对每个特征进行打分并筛选特征，聚焦于数据本身的特点。优点是计算快速、不依赖于具体模型，但缺点是选择的统计指标可能不适用于特定模型，且未考虑特征间的相互关系。

\textbf{（2）} 包裹式方法：使用模型来筛选特征，通过增加或删除特征，在验证集上测试模型准确率，寻找最优特征子集。优点是准确性较高，但缺点是计算开销大、容易过拟合。

\textbf{（3）} 嵌入式方法：将特征选择嵌入到模型构建过程中，利用模型的特性进行特征选择，如\(LASSO\)和树模型等。优点是准确率较高，计算复杂度介于过滤式和包裹式方法之间，但只有部分模型具备这个功能。



%%%%%%%%%%%%%%%%%%%%%%%%%%%%%%%%%%%%%%%%%%%%%%%%%%%%%%%%%%%%% 

\subsection{超参数寻优}

机器学习模型参数众多，参数选择不恰当，就会出现欠拟合或者过拟合的问题。为了提高模型的精度，同时提升模型的泛化能力，调参过程不可缺少。而在选择超参数的时候，有两个途径，一个是凭经验微调，另一个就是选择不同大小的参数，带入模型中，挑选表现最好的参数。

本文主要选择网格搜索，网格搜索是一种重要调参手段，即穷举搜索：在所有候选的参数选择中，通过循环遍历，尝试每一种可能性，表现最好的参数就是最终的结果。其原理就像是在数组里找到最大值，网格搜索可以保证在指定的参数范围内找到精度最高的参数。

%%%%%%%%%%%%%%%%%%%%%%%%%%%%%%%%%%%%%%%%%%%%%%%%%%%%%%%%%%%%% 

\subsection{K折交叉验证}

\(K\)折交叉验证\textsuperscript{\cite{triba2015pls}}（\(K\)-\(fold Cross Validation\)）将数据集 \( D \) 划分成 \( K \) 份互斥数据集，满足 \( D = D_1 \cup ... \cup D_k \)，一般是平均分配使每份数据量接近并且数据分布尽可能一致。每次用一份数据测试，其余 \( K - 1 \) 份数据训练，需要迭代 \( K \) 轮得到 \( K \) 个模型；最后再将 \( K \) 份测试结果汇总到一起评估一个离线指标。

\[ cv\_score = \frac{1}{K} \sum_{k=1}^{K} L(P_k, Y_k) \] 

K折交叉验证的稳定性与 \( K \) 取值有很大关系。 \( K \) 值太小实验稳定性依然偏低， \( K \) 值太大又可能导致实验成本高， \( K \) 最常用的取值是5和10。K折交叉验证能够更好地避免过拟合和欠拟合，得到的结论也更有说服力。

%%%%%%%%%%%%%%%%%%%%%%%%%%%%%%%%%%%%%%%%%%%%%%%%%%%%%%%%%%%%% 

\subsection{模型评估}

评估指标用于反映模型效果。在预测问题中，要评估模型的效果，就需要将模型预测结果 \( f(X) \) 和真实标注 \( Y \) 进行比较。
评估指标定义为 \( f(X) \) 和 \( Y \) 的函数：
\[ \text{score} = \text{metric}(f(X), Y) \]

模型的好坏是相对的，在对比不同的模型效果时，使用不同评估指标往往会导致不同的结论。

\subsubsection{准确率和召回率}

精确率和召回率多用于二分类问题，可结合混淆矩阵介绍，如下表所示。

\begin{table}[htbp]
  \centering
  \resizebox{0.5\textwidth}{!}{ 
    \begin{tabular}{@{}llll@{}}
      \toprule
                         &      & 预测结果                   &                        \\ 
      \midrule
                         &      & 正（P）                   & 负（N）                   \\
      \multicolumn{1}{c}{真实结果} & 正（P） & \multicolumn{1}{c}{TP} & \multicolumn{1}{c}{FN} \\
      \multicolumn{1}{c}{}     & 负（N） & \multicolumn{1}{c}{FP} & \multicolumn{1}{c}{TN} \\ 
      \bottomrule
    \end{tabular}
  }
  \caption{预测结果}
  \label{tab:example}
\end{table}

其中，\( TP \)（真正，\(True Positive\)）表示真实结果为正例，预测结果也是正例；\( FP \)（假正，\(False Positive\) ）表示真实结果为负例，预测结果却是正例；\( TN \)（真负，\(True Negative\)）表示真实结果为正例，预测结果却是负例；\( FN \)（假负，\(False Negative\)）表示真实结果为负例，预测结果也是负例。显然，\( TP+FP+FN+TN= \)样本总数。
精确率 \( P \) 和召回率 \( R \) 的定义为：
\[ \text{精确率}(P) = \frac{TP}{TP + FP} \quad \text{召回率}(R) = \frac{TP}{TP + FN} \]

理想情况下，精确率和召回率两者都越高越好。然而事实上这两者在某些情况下是矛盾的：精确率高时，召回率低；而精确率低时，召回率高。

\subsubsection{F1-score}

\( F_1 \) 值是精确率和召回率的调和平均值：
\[ \frac{2}{F_1} = \frac{1}{P} + \frac{1}{R} \]

\( F_1 \) 值可泛化为对精确率和召回率赋予不同权重进行加权调和：
\[ F_{\alpha} = \frac{(1+\alpha^2)\cdot P\cdot R}{\alpha^2\cdot P+R} \]

此外，准确率和错误率也是常用的评估指标。
\[ \text{准确率(\(accuracy\))} = \frac{TP + TN}{TP + FP + FN + TN} \quad \text{错误率(\(error rate)\)} = \frac{FP + FN}{TP + FP + FN + TN} \]

精确率和准确率是比较容易混淆的两个评估指标，两者是有区别的。精确率是一个二分类指标，而准确率能应用于多分类，其计算公式为：
\[ \text{准确率(\(accuracy\))} = \frac{1}{n} \sum^n_{i=1} I(f(x_i)=y_i) \]

\subsubsection{ROC 与 AUC}

在众多的机器学习模型中，很多模型输出是预测概率。而使用精确率、召回率这类指标进行模型评估时，还需要对预测概率设分类阔值，比如预测概率大于阈值为正例，反之为负例。这使得模型多了一个超参数，并且这个超参数会影响模型的泛化能力。
接收者操作特征\textsuperscript{\cite{eng2005receiver}}（\(Receiver Operating Characteristic\)，\(ROC\)）曲线不需要设定这样的阈值。\(ROC\)曲线纵坐标是真正率，横坐标是假正率，其对应的计算公式为：
\[ \text{真正率(\(TPR\))} = \frac{TP}{TP+FN} \quad \text{假正率(\(FPR\))} = \frac{FP}{TP+TN} \]

\(ROC\)曲线越靠近左上角性能越好。左上角坐标为（0,1），即 \(FPR\)=0,\(TPR\)=1，根据 \(FPR\) 和 \(TPR\) 公式可以得知，此时 \(FN\)=0,\(FP\)=0，模型对所有样本分类正确。
绘制\(ROC\)曲线很简单，首先对所有样本按预测概率排序，以每条样本的预测概率为阈值，计算对应的 \(FPR\) 和 \(TPR\)，然后用线段连接。当数据量少时，绘制的\(ROC\)曲线不平滑；当数据量大时，绘制的\(ROC\)曲线会趋于平滑。
\(AUC\)(Area Under \(ROC\) Curve）即\(ROC\)曲线下的面积，取值越大说明模型越可能将正样本排在负样本前面。\(AUC\)还有一些统计特性：\(AUC\)等于随机挑选一个正样本（P）和负样本（N）时，分类器将正样本排前面的概率；\(AUC\)和\(Wilcoxon Test of Ranks\)等价；\(AUC\)还和基尼（\(Gini\)）系数有联系，满足等式 \( Gini + I = 2 \cdot AUC \)。
\(AUC\)的计算方法有多种，从物理意义角度理解，\(AUC\)计算的是\(ROC\)曲线下的面积：
\[ AUC = \sum_{i\in(P+N)} \frac{(TPR_i+TPR_{i-1})\cdot(FPR_i-FPR_{i-1})}{2} \]

\(AUC\)计算主要与排序有关，所以它对排序敏感，而对预测分数没那么敏感。

%%%%%%%%%%%%%%%%%%%%%%%%%%%%%%%%%%%%%%%%%%%%%%%%%%%%%%%%%%%%% 

\subsection{模型融合}

\(Stacking\) \textsuperscript{\cite{chien2003sample}}是另一种更强大的模型融合方法，于 1992 年被 Wolpert 提出，其基本思路是，通过一个模型来融合若干单模型的预测结果，目的是降低单模型的泛化误差。在这里，这些单模型被称为一级模型，\(Stacking\)融合模型被称为二级模型或元模型。

\(Stacking\)先从初始的训练集训练出若干单模型，然后把单模型的输出结果作为样本特征进行整合，并把原始样本标记作为新数据样本标记，生成新的训练集。再根据新训练集训练一个新的模型，最后用新的模型对样本进行预测。\(Stacking\)融合模型本质上是一种分层结构，每一层都是若干模型。

简单起见，这里先讨论二级\(Stacking\)。同时，我们假设所有的单模型都是使用不同的学习算法产生，即异质模型（当然，\(Stacking\)的一级模型也可以是同质模型）。

在\(Stacking\)的模型训练阶段，二级模型的训练集是利用一级模型产生的。如果直接使用一级模型对初始的训练集样本进行预测来产生二级训练集，这样会有极大的过拟合的风险。因此，一般是用训练一级模型未使用的样本来产生二级模型的训练集。交叉验证法或留一法是比较常用的方法。

下面以K折交叉验证为例，来介绍二级模型的训练集是如何生成的。

输入：训练集 \( D = \{ ( x_1 , y_1 ) , ( x_2 , y_2 ) , . . . , ( x_m , y_m ) \} \)；单模型学习算法 \( \xi_1 , \xi_2 , . . . , \xi_T \)；融合模型学习算法 \( \xi \)。

\newpage
过程分析：

\begin{lstlisting}[language=Python]
for t=1,2,...,T do
    h_t=\xi_t(D)
end for
D^{'}=\phi\quad\quad //基于原始数据训练个单模型
for i=1,2,...,m do
    for t=1,2,...,T do
        z_{it}=h_t(x_i)\quad\quad //单模型输出作为样本新特征
    end for
    D^{'}=D^{'}\cup\left((z_{i1},z_{i2},...,z_{iT}),y_i\right)\quad\quad //单模型输出和样本标记构建新训练集
end for
h^{'}=\xi(D^{'})\quad\quad //在新训练集上训练二级模型
\end{lstlisting}

输出：

\( H(x) = h^{'}( h_1(x) , h_2(x) , . . . , h_T(x) ) \)  // 二级模型输出作为\(Stacking\)输出

\hspace{1em}

我们把初始训练集 \( D \) 随机划分成 \( k \) 个大小相似的集合 \( D_1 , D_2 , . . . , D_k \)。令 \( D_j \) 和 \( \bar{D}_j \) （它等于 \( D \ D_j \) ）分别表示第 \( j \) 折的测试集和训练集。给定 \( T \) 个初级学习算法，初级学习器 \( h_t^{(j)} \) 通过使用第 \( t \) 个学习算法而得。对于 \( \bar{D}_j \) 中的每个样本 \( x_i \)，定义第 \( t \) 个模型的预测结果为 \( z_{jt} = h_t^{(j)}(x_i) \)，那么由样本 \( x_i \) 产生的二级模型训练样本特征为 \( z_i = ( z_{i1} , z_{i2} , . . . , z_{iT} ) \)，样本标记还是原始的样本标记 \( y_i \)。于是，在经过 \( k \times T \) 次模型训练和预测后，得到二级训练集 \( D' = \{ ( z_i , y_i ) \} i = 1 m \)，然后 \( D' \) 用来训练二级模型。

此处二级模型 \( h' \) 是 \( ( z_1 , z_2 , . . . , z_T ) \) 关于 \( y \) 的函数。

%%%%%%%%%%%%%%%%%%%%%%%%%%%%%%%%%%%%%%%%%%%%%%%%%%%%%%%%%%%%% 

\newpage
\section{问题一}

\subsection{问题分析}

对于问题一，本团队对生产线中各装置故障数据进行了深入分析。附件 1 中的数据集涵盖了多个装置的故障情况，每个装置都具有多个特征，如物料推送状态、填装检测数、加盖数等。本团队将这些特征与装置的故障情况进行了综合分析，通过数据探索性分析，发现装置故障之间的相关性以及故障发生的潜在模式。其中，特别注意到数据集存在着类别不平衡的情况，即某些装置的故障样本数量较少，这可能会对模型的训练和泛化能力产生负面影响。因此，本团队在模型构建过程中决策分类器（\(DecisionTreeClassifier\)）来解决这一问题，以确保模型的准确性和稳定性。

\subsection{模型建立}

\subsubsection{数据处理}

\vspace{1em}

\textbf{数据准备和预处理}
\vspace{1em}

对于本题首先进行数据准备和预处理，对每个文件中的目标变量进行了二元分类处理，将非零值置为1，以适应分类模型的要求。之后分析数据并将数据可视化，将变量进行统计分析，获得各类样本数量的分布情况。

\begin{figure}[htbp]
  \centering
  \includegraphics[width=0.75\textwidth]{image_16.png}
  \caption{分类后数据}
  \label{fig:example_16}  %用于文中的图片引用
\end{figure}

\newpage
\vspace{1em}
\textbf{特征工程处理}
\vspace{1em}

在建立模型中，本团队首先对原始数据进行特征工程处理，其中包括特征选择，数据拆分以及数据标准化处理。特征选择确保模型训练所使用的特征集合与任务密切相关；数据拆分确保在训练模型之前有足够的数据用于评估模型的性能；由于本题的各个特征的数值范围差异较大，为了确保各个特征模型的影响程度相当，将数据进行标准化处理。为了处理数据集中存在的类别不平衡的情况，本题使用过采样（\(Oversampling\)）和欠采样（\(Undersampling\)）进行数据的处理。特别的是本团队对比了不同采样技术对模型性能的影响，选择最适合特定数据集和模型的技术，从而改善模型性能，减少过拟合风险，并提高对少数类样本的识别能力和召回率，使得模型更加可靠和有效。

\subsubsection{机器学习模型}

\paragraph{决策树}决策树是一种用于从数据库中导出一组规则的机器学习算法，其具有很好的可解释性，用于生成透明的决策规则。它从根节点开始，逐渐创建分支和子节点，通过对样本特征进行划分，最终生成叶节点。对于本题本团队采取了六种机器学习算法：分类与回归树模型、\(XGBoost\)回归模型、决策树回归模型、随机森林回归模型、\(Adaboost\)回归模型和\(GBDT\)回归模型。这六种模型使用的是 \(Python\) 3.6，并在\(PyCharm\) 2023.3.4环境下实现。

\paragraph{分类与回归树}\(CART\)分类与回归树（\(CART\)）是一种被广泛采用的决策树模型，其核心是应用递归划分规则来促进二叉决策树的发展。常用的规则是基尼系数（\(Gini\) \(Diversity\) \(Index\), \(GDI\)）。

\paragraph{随机森林}随机森林是一种集成模型，它利用多个决策树。RF采用了一种\(Bagging\)算法，也称为自助采样，它随机选择训练数据集的子集来训练单个决策树。自助采样过程将训练数据集随机分为“包内”子集，用于训练决策树，以及排除在培训过程之外的“包外”子集。典型的样本大小比通常设置为2:1。

\paragraph{自适应提升}自适应提升是一种提升模型，它采用自适应重抽样技术训练样本的选择。\(Adaboost\)可以利用先前训练过的决策树的信息，依次构建多个决策树。在每次迭代中，数据集被分配权重，加大对先前错误分类样本的关注，错误分类的样本比正确分类的样本更容易被优先选择。这一策略使新的决策树能够在处理新数据集时表现更加出色。

\paragraph{极端梯度提升}极端梯度提升代表了\(GBM\)的增强版本，其中引入了二阶导数将损失函数最小化，从而产生更精确的决策树。在\(XGBoost\)框架内，收缩和列抽样策略在减轻模型偏差和方差方面发挥了关键作用。具体来说，收缩减弱了个别树的影响，从而降低了偏差，而列抽样引入随机性到模型中，有效减小了方差。此外，\(XGBoost\)还引入了正则化项来抵消过拟合倾向。

\subsection{问题求解}
\paragraph{模型训练}
\subparagraph{不同模型训练}
对于问题一，首先在不处理类别不平衡的情况下使用上述五种模型分别进行训练，其中训练集占70\%，测试集占30\%。

\begin{figure}[htbp]
  \centering
  \includegraphics[width=0.85\textwidth]{image_17.png}
  \caption{数据处理前模型性能比较}
  \label{fig:example_17}  %用于文中的图片引用
\end{figure}

使用原始数据进行模型训练得到的结果出现过拟合的问题，本团队分别使用过采样和欠采样对数据进行处理得到结果。


\begin{table}[h]
\centering
\resizebox{0.75\textwidth}{!}{%
\begin{tabular}{llllll}
\hline
Model    & Accuracy & Precision & Recall & Specificity & F1 score \\ \hline
CART     & 0.6693   & 0.7009    & 0.6705 & 0.6678      & 0.6854   \\
RF       & 0.7624   & 0.7644    & 0.7713 & 0.7532      & 0.7678   \\
AdaBoost & 0.7329   & 0.7221    & 0.7492 & 0.7169      & 0.7354   \\
GBM      & 0.7655   & 0.7674    & 0.7744 & 0.7563      & 0.7709   \\
XGBoost  & 0.7842   & 0.7795    & 0.7963 & 0.7719      & 0.7878   \\
LightGBM & 0.7857   & 0.7885    & 0.7933 & 0.7778      & 0.7909   \\ \hline
\end{tabular}%
}
\caption{数据处理后模型性能比较}
\end{table}

\begin{figure}[htbp]
  \centering
  \includegraphics[width=0.85\textwidth]{image_18.png}
  \caption{数据处理后模型性能比较}
  \label{fig:example_1}  %用于文中的图片引用
\end{figure}



\newpage
\subsection{总结}

\(LightGBM\)模型在\(Accuracy\)、\(Precision\)、\(Specificity\)和\(F1 score\)方面的数值最高，这体现了该模型在处理故障数据的卓越性能。同时，在\(Recall\)方面，\(LightGBM\)模型也排名第二。\(XGBoost\)模型在所有性能指标上的表现紧随其后，排名第二。然而，\(CART\)模型在所有性能指标上都记录了最低的数值，这表明该模型的建模性能相对较差。其次，\(AdaBoost\)模型的预测性能也相对较差。

考虑到各种性能指标，可以得出结论：在滑坡易发性制图中，基于提升算法的三种集成学习模型\(LightGBM\)、\(XGBoost\)和\(GBM\)拥有卓越性能。相反\(XGBoost\)模型在性能方面超越了\(LightGBM\)模型。实际上，选择模型类型取决于特定的研究区域和变量，模型的适用性和性能可能会相应变化。一些研究表明，与RF相比，提升算法可以显著提高模型的\(Accuracy\)。\(XGBoost\)模型在缓解过拟合问题方面表现出明显优势，而\(LightGBM\)模型则依赖直方图进行快速计算和降低内存消耗，使其非常适合处理大规模数据集。在特定研究场景中，有限的样本大小可能是导致其次优性能的关键因素。对于本题我们选择了使用\(LightGBM\)模型。

不同故障类型模型训练在本题中，我们分别针对加盖装置定位故障5001、加盖装置加盖故障5002、拧盖装置定位故障6001、拧盖装置拧盖故障6002四种不同类型的故障训练了相应的\(LightGBM\)模型，并评估了模型的性能。

通过训练模型，可以对生产过程中可能出现的不同故障类型进行识别和预测。一旦故障发生，系统可以及时发出警报或采取相应措施，以减少生产线停机时间和生产成本。及时发现和解决生产过程中的故障可以提高生产效率和产品质量。通过建立模型，可以发现故障出现的模式和规律，从而优化生产流程，提高生产效率。基于历史数据和机器学习模型，可以实现设备的智能维护，提前预测设备的故障并进行维修，从而减少意外停机和维修成本。通过不断监测和分析生产过程中的数据，可以发现潜在的问题和改进空间，从而持续改进生产流程和设备性能。

%%%%%%%%%%%%%%%%%%%%%%%%%%%%%%%%%%%%%%%%%%%%%%%%%%%%%%%%%%%%% 

\newpage
\section{问题二}

在问题一中建立了六种模型模型，分别对比了六种模型的评分以及使用过采样和欠采样得分，最终选择LightGBM模型。根据该模型对附件2中的数据进行分析判断得出故障报警的日期，开始时间和持续时长见下表。

\begin{table}[h]
\centering
\begin{adjustbox}{width=0.65\textwidth}
\begin{tabular}{llll}
\toprule
生产线编号 & 故障总次数 & 最长持续时间 & 最短持续时间 \\ \midrule
1001  & 32 & 167    & 1      \\
2001  & 33          & 9      & 1      \\
4001  & 41          & 187    & 1      \\
4002  & 150         & 172    & 1      \\
4003  & 2825        & 54     & 1      \\
5001  & 2825        & 18     & 1      \\
5002  & 2825        & 13     & 1      \\
6001  & 2825        & 143    & 1      \\
6002  & 3699        & 72     & 1      \\ \bottomrule
\end{tabular}
\end{adjustbox}
\caption{\(LightGBM\)的数据分析结果}
\end{table}

%%%%%%%%%%%%%%%%%%%%%%%%%%%%%%%%%%%%%%%%%%%%%%%%%%%%%%%%%%%%% 

\newpage
\section{问题三}

\subsection{问题分析}

\subsubsection{数据汇总}

要分析产品的产量、合格率与生产线、操作人员等因素之间的关系，本团队认为首先需要将生产线的运行记录数据汇总成一条数据。通过设计一些能够代表生产线在一年内的整体运行情况汇总变量，提高后序数据分析逻辑性和准确性，下面是一些使用到的汇总变量、计算公式及其原理：

\textbf{（1）} 总运行时间：计算生产线在一年内运行的总时间。

\begin{itemize}
    \item 公式：总运行时间 = ∑（结束时间 - 开始时间）
    \item 原理：累加每个生产周期的运行时间。
\end{itemize}

\textbf{（2）} 总生产数量：计算一年内生产线生产的总产品数量。

\begin{itemize}
    \item 公式：总生产数量 = ∑合格数
    \item 原理：累加每天记录的合格产品数量。
\end{itemize}

\textbf{（3）} 总不合格数量：计算一年内生产线生产的不合格产品总数量。

\begin{itemize}
    \item 公式：总不合格数量 = ∑不合格数
    \item 原理：累加每天记录的不合格产品数量。
\end{itemize}

\textbf{（4）} 平均生产效率：计算生产线的平均生产效率，即每小时生产的合格产品数量。

\begin{itemize}
    \item 公式：平均生产效率 = 总生产数量 / 总运行时间
    \item 原理：用总生产数量除以总运行时间得到每小时的生产效率。
\end{itemize}

\textbf{（5）} 设备综合故障率：计算所有设备故障的综合故障率。

\begin{itemize}
    \item 公式：设备综合故障率 = （物料推送装置故障1001 + 物料检测装置故障2001 + ... + 拧盖装置拧盖故障6002） / 总天数 * 100%`
    \item 原理：累加所有设备故障的发生次数，然后除以一年的总天数，得到设备综合故障率。
\end{itemize}

\textbf{（6）} 物料推送效率：计算物料推送气缸的平均推送效率。

\begin{itemize}
    \item 公式：物料推送效率 = （物料推送数 / (物料推送数 + 物料待抓取数)） * 100%
    \item 原理：用成功推送的物料数量除以（成功推送的物料数量加上待抓取的物料数量），得到推送效率。
\end{itemize}

\textbf{（7）} 填装效率：计算填装过程的效率。

\begin{itemize}
    \item 公式：填装效率 = （填装数 / 物料推送数） * 100%
    \item 原理：用完成填装的次数除以成功推送的物料数量，得到填装效率。
\end{itemize}

\textbf{（8）} 加盖效率：计算加盖过程的效率。

\begin{itemize}
    \item 公式：加盖效率 = （加盖数 / 填装数） * 100%
    \item 原理：用完成加盖的次数除以完成填装的次数，得到加盖效率。
\end{itemize}

\textbf{（9）} 拧盖效率：计算拧盖过程的效率。

\begin{itemize}
    \item 公式：拧盖效率 = （拧盖数 / 加盖数） * 100%
    \item 原理：用完成拧盖的次数除以完成加盖的次数，得到拧盖效率。
\end{itemize}

\textbf{（10）} 生产总周期数：计算生产线一年内的总生产周期数。

\begin{itemize}
    \item 公式：生产总周期数 = ∑填装数
    \item 原理：累加每天记录的填装次数，得到总生产周期数。
\end{itemize}

\textbf{（11）} 合格率：计算生产线的合格率。

\begin{itemize}
    \item 公式：合格率 = （总生产数量 / (总生产数量 + 总不合格数量)） * 100%
    \item 原理：用总生产数量除以（总生产数量加上总不合格数量），得到合格率。
\end{itemize}

这些汇总变量可以帮助本团队从不同角度评估生产线的运行状况，包括生产效率、产品质量、设备可靠性等。在后序建模过程中中，本团队根据需要选择适当的汇总变量进行分析。

\subsubsection{数据结构}

本团队在对数据进行汇总之后，进一步对数据进行处理，处理流程如下：

\textbf{(1)} 数据处理和读取

\begin{itemize}
    \item 遍历读取了附件3中的所有\(CSV\)文件，并将数据存储在\(DataFrame\)中。
    \item 日期列进行了处理，将日期列转换为日期格式，并提取唯一日期值以遍历每个日期。
\end{itemize}

\textbf{(2)} 统计分析

\begin{itemize}
    \item 对每个日期进行统计分析，计算了各项生产指标的数值。
    \item 包括总运行时间、总生产数量、总不合格数量、平均生产效率、设备综合故障率、物料推送效率、填装效率、加盖效率、拧盖效率、生产总周期数和合格率等指标。
\end{itemize}

\textbf{(3)} 数据整合

\begin{itemize}
    \item 将每个日期计算得到的指标数据整合到一个\(DataFrame\)中。
    \item 在整合数据时，本团队还将操作人员编号、工龄和生产线编号等信息与指标数据对应起来。
\end{itemize}

\textbf{(4)} 迭代处理

\begin{itemize}
    \item 通过迭代遍历每个\(CSV\)文件和每个日期，实现了对全部数据的统计分析和整合。
\end{itemize}

这段代码的作用是对附件 3 中的数据进行了深入分析，计算了各种生产指标如：总运行时间、总生产数量、总不合格数量、平均生产效率、设备综合故障率物料推送效率等生产指标以及不同工龄对应的合格率以及对应的总生产数量
，并将结果整合在一个\(DataFrame\)中，为后续对产品产量、合格率与生产线、操作人员等因素之间的关系进行详细分析提供了数据基础。

\subsubsection{模型选择}

本团队使用了多种回归模型对规整后数据集进行了建模和评估，并且使用了\(SHAP\)（\(SHapley Additive exPlanations\)）值来解释模型的预测结果。

\textbf{（1）}线性回归模型：

使用线性回归模型对数据进行拟合和预测。线性回归模型简单易懂，对于线性关系较强的数据适用性较好，因此被选用作为基准模型。

\textbf{（2）}决策树回归模型：

决策树模型能够处理非线性关系，并且具有较好的解释性，容易理解模型的决策过程，因此被用于对复杂数据进行建模。

\textbf{（3）}随机森林回归模型：

随机森林是一种集成学习算法，通过集成多个决策树来提高模型的预测性能和稳定性，对于高维度数据和大规模数据集有较好的适应性。

\textbf{（4）}支持向量回归模型：

支持向量机能够处理高维度数据和非线性关系，具有较强的泛化能力和稳健性，适用于复杂的数据建模任务。

\textbf{（5）}神经网络回归模型：

神经网络模型是一种强大的非线性回归模型，能够学习复杂的数据模式和关系，对于高度非线性的数据建模具有较好的效果。

本团队选择这些模型进行数据处理是基于它们的特点和适用范围，通过对比不同模型的评价指标，可以综合考虑模型的预测性能和泛化能力，从而选取最适合数据特征的模型进行建模和预测。同时，使用\(SHAP\)值对模型进行解释，可以帮助理解模型对于不同特征的重要性，从而进一步优化模型和数据特征的选择。

%%%%%%%%%%%%%%%%%%%%%%%%%%%%%%%%%%%%%%%%%%%%%%%%%%%%%%%%%%%%% 
\subsection{模型建立}

这段代码是一个典型的机器学习模型建立与评估的流程。首先，我们导入了需要的模块和库，包括了不同类型的回归模型和评价指标的计算方法。接着，我们按照以下步骤逐一建立了不同的回归模型，并对其进行评估。

\textbf{（1）} 线性回归模型：

\begin{itemize}
    \item 线性回归是最简单直观的回归模型之一，它假设自变量与因变量之间存在线性关系。
    
    \item 自变量X：{ 总运行时间，总生产数量，总不合格数量，平均生产效率，设备综合故障率，物料推送效率，填装效率，加盖效率，拧盖效率，生产总周期数，合格率 }；
    
    \item 因变量Y：{ 工龄 }
\end{itemize}

这段代码中，本团队使用了\(LinearRegression\)类来构建线性回归模型，并通过计算均方误差\(MSE\) 和决定系数来评估模型的拟合效果和预测准确性

\textbf{（2）} 决策树回归模型：

\begin{itemize}
    \item 决策树是一种树形结构的分类模型，通过对数据集进行多次分割，将数据划分到不同的叶子节点中。
    
    \item 决策树回归模型通过划分自变量的取值范围，将数据划分为不同的区域，并在每个区域内使用区域内的平均值来进行预测。
    
\end{itemize}

本团队使用\(DecisionTreeRegressor\)类来构建决策树回归模型，并选择与线性回归模型相同的系数来评估模型性能。


\textbf{（3）} 随机森林回归模型：

\begin{itemize}

    \item 随机森林回归模型在决策树的基础上，进一步引入了随机性，从而提高了模型的泛化能力和预测准确性。
    
\end{itemize}

在代码中，本团队使用\(RandomForestRegressor\)类来构建随机森林回归模型，并对模型进行评估。

\textbf{（4）} 支持向量回归模型：

\begin{itemize}
    \item 支持向量机是一种强大的分类与回归方法，它通过寻找最佳的超平面来进行分类或回归。
    
    \item 支持向量机回归模型通过找到训练数据中的支持向量，并以它们为基础构建回归函数。
    
\end{itemize}

本团队使用\(SVR\)类来构建支持向量机回归模型，并对模型性能进行评估。

\textbf{（5）} 神经网络回归回归模型：

 本团队使用\(MLPRegressor\)类来构建神经网络回归模型，并对模型性能进行评估。


通过以上步骤，本团队可以得到不同回归模型的预测结果，并通过评价指标对模型进行评估，从而选择最合适的模型用于实际预测任务。

%%%%%%%%%%%%%%%%%%%%%%%%%%%%%%%%%%%%%%%%%%%%%%%%%%%%%%%%%%%%% 

\subsection{问题求解}

\subsubsection{求解建模}

本团队首先导入了所需的模块和库，包括不同类型的回归模型和评价指标的计算方法。接着，按照以下步骤逐一建立了不同的回归模型：

\textbf{(1)} 使用线性回归模型，它是最简单直观的回归模型之一，通过对数据拟合直线来预测连续型变量的值。

\textbf{(2)} 构建了决策树回归模型，它通过划分自变量的取值范围，将数据划分为不同的区域，并在每个区域内使用区域内的平均值来进行预测。

\textbf{(3)} 使用随机森林回归模型，它通过构建多个决策树，并综合它们的预测结果来进行最终预测，提高了模型的泛化能力和预测准确性。

\textbf{(4)} 构建了支持向量机回归模型，它通过找到训练数据中的支持向量，并以它们为基础构建回归函数。

\textbf{(5)} 使用神经网络回归模型，通过多层神经元的连接和激活函数的作用，对数据进行复杂的非线性建模。


对于每个模型，本团队都对其进行了评估，包括计算均方误差和决定系数等指标，以评估模型的拟合效果和预测准确性。通过这些步骤，本团队可以得到不同回归模型的预测结果，并选择最合适的模型用于实际预测任务。

\subsubsection{模型输出}

\textbf{(1)} 数据汇总
\vspace{1em}

数据汇总后有如下输出：

工人工龄及其生产率和总生产数量的数据统计如表所示：

\begin{table}[htbp]
  \centering
  \begin{adjustbox}{width=0.5\textwidth}
    \begin{tabular}{lll}
    \toprule
    工龄 & 合格率 & 总生产量 \\
    \midrule
    1  & 99.928919 & 81297.510385 \\
    2  & 99.947596 & 81328.768572 \\
    3  & 99.904792 & 79901.154734 \\
    4  & 99.903422 & 81499.468565 \\
    5  & 99.926576 & 81534.313247 \\
    6  & 99.963954 & 83035.108491 \\
    \bottomrule
    \end{tabular}%
  \end{adjustbox}
  \caption{工人年龄及其生产效率和总生产数量的统计}
  \label{tab:addlabel}%
\end{table}%

\newpage
\textbf{（2）}各模型评价结果
\vspace{1em}

由本团队采用的五种模型的输出结果实际值与预测值如图所示：

\begin{figure}[htbp]
  \centering
  \includegraphics[width=0.95\textwidth]{image_8.png}
  \caption{各模型评价结果}
  \label{fig:example_8}  %用于文中的图片引用
\end{figure}

评价结果如下表所示：

\begin{table}[htbp]
  \centering
  \begin{adjustbox}{width=0.75\textwidth}
    \begin{tabular}{lll}
    \toprule
    \textbf{} & \textbf{MSE} & \textbf{$R^2$} \\
    \midrule
    线性回归 & 2.6089198956304247 & 0.07155875600340755 \\
    决策树模型 & 0.48644871794871797 & 0.8268865772424492 \\
    随机森林模型 & 0.7240109048744705 & 0.7423448737101528 \\
    支持向量机模型 & 2.721408220857354 & 0.03152732353830823 \\
    神经网络模型 & 48641.26341272193 & -17309.058153993567 \\
    \bottomrule
    \end{tabular}%
  \end{adjustbox}
  \caption{各模型评价结果}
  \label{tab:addlabel}%
\end{table}%

这个表格展示了本团队使用不同回归模型对数据集进行预测时的评价指标，包括均方误差（\(MSE\)）和决定系数（$R^2$）。下面是每个模型的数据意义和结论：

\begin{itemize}
    \item 线性回归模型：MSE为2.61，$R^2$为0.07。线性回归模型的预测效果较差，MSE较高，说明模型的预测值与实际值之间的偏差较大，$R^2$接近于0，说明模型对数据的拟合程度较低。

    \item 决策树模型：MSE为0.49，$R^2$为0.83。决策树模型表现良好，MSE较低，$R^2$接近于1，说明模型对数据的拟合程度较高，预测精度较高。

    \item 随机森林模型：MSE为0.72，$R^2$为0.74。随机森林模型的预测效果较好，MSE较低，$R^2$接近于1，说明模型对数据的拟合程度较高，预测精度较高。

    \item 支持向量机模型：MSE为2.72，$R^2$为0.03。支持向量机模型的预测效果较差，MSE较高，$R^2$接近于0，说明模型对数据的拟合程度较低。

    \item 神经网络模型：MSE为48641.26，$R^2$为-17309.06。神经网络模型的预测效果非常差，MSE极高且$R^2$为负值，说明模型对数据的拟合程度非常低，预测精度极差。
    
\end{itemize}

综上所述，决策树模型和随机森林模型在这个数据集上表现最好，预测精度较高，而神经网络模型的表现最差，预测效果极差。

\vspace{1em}

\textbf{（3）}SHAP

\vspace{1em}

\(SHAP\)（\(SHapley Additive exPlanations\)）是一种用于解释机器学习模型预测结果的方法。它基于合作博弈论中的 \(Shapley\) 值理论，通过计算每个特征对于模型预测结果的贡献度来解释模型的预测过程。值提供了对于每个特征的影响程度，使得我们能够更深入地理解模型是如何基于输入特征来做出预测的。

在给定数据集和训练好的模型后，\(SHAP\)库可以计算每个特征对于每个样本的\(SHAP\)值。这些\(SHAP\) 值可以用于解释模型对于单个样本的预测，也可以用于对整个数据集的预测进行总体解释。\(SHAP\)值的正负表示特征对于增加还是减少预测结果的影响，而数值的大小表示影响的程度。

通过可视化\(SHAP\)值，我们可以直观地理解模型对于每个特征的重要性，以及特征值的不同对于模型预测的影响。这有助于我们对模型的预测结果进行解释，并且可以帮助我们发现数据中的模式和规律。

对各特征\(SHAP\)值可视化，如表所示：

\begin{figure}[htbp]
  \centering
  \includegraphics[width=0.75\textwidth]{image_9.png}
  \caption{\(SHAP\)值可视化}
  \label{fig:example_9}  %用于文中的图片引用
\end{figure}

下面是每个特征的含义及其对模型预测的贡献：

\begin{itemize}
    \item 平均生产效率：对模型预测结果的贡献最大，SHAP值为0.431，说明平均生产效率对于模型的预测结果有较大的影响。

    \item 加盖效率和填装效率：对模型预测结果的贡献较大，SHAP值分别为0.248和0.247，说明加盖效率和填装效率也对模型的预测结果有一定的影响。
    
    \item 拧盖效率、生产总周期数、物料推送效率和总运行时间：对模型预测结果的贡献相对较小，但仍然有一定的影响。

    \item 设备综合故障率、总生产数量、总不合格数量和合格率：对模型预测结果的贡献较小，SHAP值较低，说明这些特征对于模型的预测结果影响较小。
    
\end{itemize}

综上所述，本团队发现平均生产效率、加盖效率和填装效率是对模型预测结果影响较大的关键特征，而设备综合故障率、总生产数量、总不合格数量和合格率对模型预测结果的影响相对较小。

\begin{figure}[htbp]
  \centering
  \includegraphics[width=0.8\textwidth]{image_10.png}
  \caption{\(SHAP\)摘要图}
  \label{fig:example_10}  %用于文中的图片引用
\end{figure}

%%%%%%%%%%%%%%%%%%%%%%%%%%%%%%%%%%%%%%%%%%%%%%%%%%%%%%%%%%%%% 

\subsection{总结}


\textbf{针对工龄的多独立样本\(Kruskal\)-\(Wallis\)检验}

\subsubsection{分析流程}

\textbf{数据源}

\vspace{1em}
部分数据截图如图所示：

\begin{figure}[htbp]
  \centering
  \includegraphics[width=0.95\textwidth]{image_11.png}
  \caption{\(SHAP\)摘要图}
  \label{fig:example_11}  %用于文中的图片引用
\end{figure}

\textbf{算法配置}

\vspace{1em}
\begin{itemize}
    \item 算法：多独立样本\(Kruskal\)-\(Wallis\)检验  
    \item 分组变量：{ 工龄 }；变量Y：{ 总生产数量，合格率，平均生产效率，总不合格数量 }  
\end{itemize}

\vspace{1em}
\textbf{结果分析}
\vspace{1em}

多独立样本\(Kruskal\)-\(Wallis\)检验用于分析多分组数据是否存在显著性差异：

\begin{itemize}

    \item 基于变量\textbf{总生产数量}，检验结果P值为0.000<0.05，因此统计结果显著，说明不同工龄在总生产数量上存在显著差异；其差异幅度Cohen's f值为：0.025，极小程度差异。 
    
    \item 基于变量\textbf{合格率}，检验结果P值为0.001<0.05，因此统计结果显著，说明不同工龄在合格率上存在显著差异；其差异幅度Cohen's f值为：0.004，极小程度差异。 

    \item 基于变量\textbf{平均生产效率}，检验结果P值为0.000<0.05，因此统计结果显著，说明不同工龄在平均生产效率上存在显著差异；其差异幅度Cohen's f值为：0.025，极小程度差异。 

    \item 基于变量\textbf{总不合格数量}，检验结果P值为0.001<0.05，因此统计结果显著，说明不同工龄在总不合格数量上存在显著差异；其差异幅度Cohen's f值为：0.004，极小程度差异。
    
\end{itemize}

\newpage
\vspace{1em}
\textbf{详细结论}
\vspace{1em}

输出结果：正态性检验直方图

\textbf{(1)} 总生产量

\begin{figure}[htbp]
  \centering
  \includegraphics[width=0.75\textwidth]{image_12.jpg}
  \caption{总生产数量}
  \label{fig:example_12}  %用于文中的图片引用
\end{figure}

\textbf{(2)} 合格率

\begin{figure}[htbp]
  \centering
  \includegraphics[width=0.75\textwidth]{image_13.jpg}
  \caption{合格率}
  \label{fig:example_13}  %用于文中的图片引用
\end{figure}

\newpage

\textbf{(3)} 平均生产效率

\begin{figure}[htbp]
  \centering
  \includegraphics[width=0.75\textwidth]{image_14.jpg}
  \caption{平均生产效率}
  \label{fig:example_14}  %用于文中的图片引用
\end{figure}

\textbf{(4)} 总不合格数量

\begin{figure}[htbp]
  \centering
  \includegraphics[width=0.75\textwidth]{image_15.jpg}
  \caption{总不合格数量}
  \label{fig:example_15}  %用于文中的图片引用
\end{figure}

\textbf{（5）} \(Kruskal\)-\(Wallis\)检验分析结果表

\begin{table}[htbp]
  \centering
  \resizebox{0.95\textwidth}{!}{
    \begin{tabular}{cccccccc}
    \toprule
    \textbf{分析项} & \textbf{分组变量} & \textbf{样本量} & \textbf{中位数} & \textbf{标准差} & \textbf{统计量} & \textbf{P} & \textbf{Cohen's f值} \\ \midrule
    总生产数量        & 1             & 520          & 21039350.5   & 479985.468   & 385.27       & 0.000***   & 0.025               \\
    5            & 520           & 21081588     & 423861.395   &              &              &            &                     \\
    6            & 260           & 21621056.5   & 234415.468   &              &              &            &                     \\
    4            & 520           & 21065010     & 440415.567   &              &              &            &                     \\
    2            & 520           & 21056912     & 485797.915   &              &              &            &                     \\
    3            & 260           & 20810334.5   & 233733.449   &              &              &            &                     \\
    总计           & 2600          & 21061091.5   & 460956.105   &              &              &            &                     \\ \midrule
    合格率          & 1             & 520          & 100          & 0.26         & 20.528       & 0.001***   & 0.004               \\ \midrule
    5            & 520           & 100          & 0.261        &              &              &            &                     \\
    6            & 260           & 100          & 0.167        &              &              &            &                     \\
    4            & 520           & 100          & 0.293        &              &              &            &                     \\
    2            & 520           & 100          & 0.206        &              &              &            &                     \\
    3            & 260           & 100          & 0.284        &              &              &            &                     \\
    总计           & 2600          & 100          & 0.253        &              &              &            &                     \\ \midrule
    平均生产效率       & 1             & 520          & 724.089      & 16.156       & 397.706      & 0.000***   & 0.025               \\ \midrule
    5            & 520           & 724.792      & 14.307       &              &              &            &                     \\
    6            & 260           & 746.991      & 6.999        &              &              &            &                     \\
    4            & 520           & 724.589      & 14.692       &              &              &            &                     \\
    2            & 520           & 724.402      & 16.381       &              &              &            &                     \\
    3            & 260           & 718.823      & 7.147        &              &              &            &                     \\
    总计           & 2600          & 724.523      & 15.496       &              &              &            &                     \\ \midrule
    总不合格数量       & 1             & 520          & 0            & 54801.363    & 20.471       & 0.001***   & 0.004               \\ \midrule
    5            & 520           & 0            & 54804.57     &              &              &            &                     \\
    6            & 260           & 0            & 36146.577    &              &              &            &                     \\
    4            & 520           & 0            & 62096.988    &              &              &            &                     \\
    2            & 520           & 0            & 43484.023    &              &              &            &                     \\
    3            & 260           & 0            & 58894.814    &              &              &            &                     \\
    总计           & 2600          & 0            & 53286.786    &              &              &            &                     \\ \midrule
    \multicolumn{8}{c}{注：***、**、*分别代表1\%、5\%、10\%的显著性水平}                                                                           \\ \bottomrule
    \end{tabular}
  }
  \caption{Kruskal-Wallis检验分析结果}
  \label{tab:kruskal-wallis}
\end{table}

\newpage
\textbf{分析：}

\begin{itemize}

    \item \(Kruskal-Wallis\)检验结果显示，基于变量总生产数量，检验结果 P 值为 0.000<0.05，因此统计结果显著，说明不同工龄在总生产数量上存在显著差异；其差异幅度 \(Cohen's f\)值为：0.025，极小程度差异。 
    
    \item \(Kruskal-Wallis\)检验结果显示，基于变量合格率，检验结果P值为0.001<0.05，因此统计结果显著，说明不同工龄在合格率上存在显著差异；其差异幅度\(Cohen's f\)值为：0.004，极小程度差异。  

    \item \(Kruskal-Wallis\)检验结果显示，基于变量平均生产效率，检验结果P值为0.000<0.05，因此统计结果显著，说明不同工龄在平均生产效率上存在显著差异；其差异幅度\(Cohen's f\)值为：0.025，极小程度差异。 

    \item \(Kruskal-Wallis\)检验结果显示，基于变量总不合格数量，检验结果P值为0.001<0.05，因此统计结果显著，说明不同工龄在总不合格数量上存在显著差异；其差异幅度\(Cohen's f\)值为：0.004，极小程度差异。  
 
\end{itemize}

%%%%%%%%%%%%%%%%%%%%%%%%%%%%%%%%%%%%%%%%%%%%%%%%%%%%%%%%%%%%% 

\newpage
\section{问题四}

\subsection{问题分析}
为制定最佳的操作人员排班方案，本团队首先根据问题三的分析结果，精确确定各个工龄段操作人员在整体操作人员中的比例。这一步骤的关键在于确保排班方案能够充分考虑到不同工龄段操作人员的数量分布，从而有效利用各个工龄段的人力资源，提高生产效率。

制定过程中，本团队利用聚类分析等机器学习模型，对操作人员的工龄、工作经验等特征进行聚类，以便更好地理解和描述不同工龄段操作人员的特征和数量分布情况。这种数据驱动的方法，使得评估不同工龄段操作人员在排班方案中的角色和贡献更加客观，从而为制定最佳排班方案提供更有针对性的参考。

根据各个工龄段操作人员的比例以及总操作人员数量，本团队采用决策树等机器学习模型，来预测和优化不同工龄段操作人员在不同班次中的分配情况。此种预测性的建模方法，有利于理解和预测不同操作人员在不同时间段的工作能力和生产效率，从而更加灵活安排他们的工作时间和班次分配。

确定了操作人员的班次分配后，本团队着重考虑操作人员的工作时间安排。为了保障操作人员的身心健康和工作效率，利用时间序列分析等机器学习模型，来分析和预测不同班次中的工作负荷和生产线的运行状态。这种数据驱动的方法使得评估和优化操作人员的工作时间安排更加精确，以最大程度地提高生产效率和保障工人的健康和安全。

此外，为了最大程度地提高操作人员的工作效率和生产线的运行稳定性，本团队利用\(SVM\)深度学习模型分析和优化操作人员的工作行为和生产过程。通过这种基于数据的深度学习方法，可以更加细致地理解和优化生产线中的各个环节和操作过程，从而提高生产效率和产品质量。


\subsection{模型建立}

\subsubsection{目标函数}

目标是最大化合格率乘以生产数量的总和。目标函数的公式可以表示为：

\[
\text{Maximize } Z = \sum_{i=1}^{N} \sum_{j=1}^{D} \sum_{k \in S} x_{ijk} \cdot E_{\text{rate},i} \cdot E_{\text{quantity},i}
\]

其中：

\begin{itemize}
    \item \( N \) 是工人总数\(num_workers\)。
    \item \( D \) 是一周的天数\(num_days\)。
    \item \(S\)是班次集合\(shifts\)，包括早、中、晚班。
    \item \( x_{ijk} \) 是决策变量，表示工人 \( i \) 在第 \( j \) 天的班次 \( k \) 是否工作（1表示工作，0表示不工作）。
    \item  \( E_{\text{rate},i} \) 是工人 \( i \) 的合格率。
    \item  \( E_{\text{quantity},i} \) 是工人 \( i \) 的每日生产数量。
    
\end{itemize}

\subsubsection{约束条件 (Constraints)}

\textbf{（1）} 工作天数约束：每个工人每周工作 5 天，休息 2 天。

\[
\sum_{j=1}^{D} \sum_{k \in S} x_{ij} = 5, \quad \forall i = 1, 2, ..., N
\]

\textbf{（2）} 班次人数约束：每个班次每天需要 10 个不同的人。

\[
\sum_{i=1}^{N} x_{ijk} = 10, \quad \forall j = 1, 2, ..., D, \quad \forall k \in S
\]

\textbf{（3）} 工龄分布比例约束：确保每个工龄等级的工人数量符合给定的比例。

\[
\sum_{i=1}^{N} \sum_{j=1}^{D} \sum_{k \in S} x_{ijk} = W_{\text{exp}}, \quad \forall \text{exp} \in \text{experience\_levels}
\]

其中 \( W_{\text{exp}} \) 是工龄等级 \( \text{exp} \) 的工人在一周内工作的总天数，根据工龄等级的比例和工人总数计算得出。

\subsubsection{模型构建}

构建的模型是一个线性规划模型，旨在优化工人排班方案以最大化生产质量。本模型使用了\(Python\)中的\(pulp\)库来构建和求解。下面对该模型进行分析：

\textbf{（1）} 目标函数：目标函数旨在最大化合格率乘以生产数量的总和。这个目标函数反映了对生产线整体质量的优化需求，使得每个班次的生产效率和质量最大化。

\textbf{（2）} 决策变量：决策变量\(x_{ij}\)表示工人i在第j天的班次安排，取值为二进制变量，表示工人是否参与了对应的班次。这些决策变量构成了整个排班方案的基础。

\textbf{（3）} 约束条件：

\begin{itemize}
    \item 每个工人每周工作 5 天，休息 2 天：通过约束每个工人在一周内的班次总和为 5 来实现。
    \item 每个班次每天需要 10 个不同的人：通过约束每个班次在每天的工人总和等于 10 来实现。
    \item 工龄分布比例约束：根据给定的工龄分布比例，约束不同工龄段工人在排班方案中的数量分布。
\end{itemize}

\textbf{（4）} 求解方法：利用 \(pulp\) 库中的\(solve\)() 方法求解该线性规划模型，得到最优的工人排班方案。

在构建模型后，通过将决策变量的值转换为班次或休息的方式，创建了一个\(DataFrame\)数据表，展示了每个工人在每天的班次安排情况。这个数据表可以进一步用于分析和可视化排班方案的具体实施情况和效果。

\begin{table}[htbp]
  \centering
  \resizebox{0.65\textwidth}{!}{ % 将表格调整为页面宽度的0.75倍
    \begin{tabular}{clcc}
      \toprule
      \multicolumn{1}{l}{工龄（年）} & 合格率                           & 总生产数量    & \multicolumn{1}{l}{人员占比要求} \\ 
      \midrule
      1                         & \multicolumn{1}{c}{99.943949} & 21555903 & 0.2                        \\
      2                         & 99.947722                     & 21145480 & 0.2                        \\
      3                         & 99.905039                     & 20774300 & 0.1                        \\
      4                         & 99.903429                     & 21189862 & 0.2                        \\
      5                         & 99.926514                     & 21198921 & 0.2                        \\
      6                         & 99.963823                     & 21589128 & 0.1                        \\ 
      \bottomrule
    \end{tabular}
  }
  \caption{排版方案具体实施情况和效果}
  \label{tab:example}
\end{table}

\subsection{问题求解}

\subsubsection{求解建模}

此问题可以建模为一个混合整数线性规划\(MILP\)问题，其中使用\(PULP\)库进行求解。

\(MILP\)的数学模型如下表示：

\[ \text{Maximize:} \quad Z = \sum_{i=1}^{N} \sum_{j=1}^{D} \sum_{k=1}^{S} (x_{i,j,k} \times \text{rate}_{i} \times \text{quantity}_{i}) \]

其中，

\begin{itemize}

    \item \( Z \) 表示总的生产质量；
    \item \( N \) 为工人总数；
    \item \( D \) 为每周的天数；
    \item \( S \) 为班次数量；
    \item \( x_{i,j,k}\) 为决策变量，表示工人 \( i \) 在第 \( j \) 天的第 \( k \) 班次的分配情况；
    \item \(\text{rate}_{i} \) 表示工人 \( i \) 的合格率；
    \item \(\text{quantity}_{i} \) 表示工人 \( i \) 的每日生产数量。
    
\end{itemize}

\vspace{12pt}
为了求解\(MILP\)问题，本团队按照以下步骤进行求解：

\textbf{(1)} 初始化问题：创建一个空的线性规划模型。

\textbf{(2)} 定义决策变量：创建二进制决策变量，表示每个工人在每天的每个班次的分配情况。

\textbf{(3)} 构建目标函数：定义目标函数，即最大化总的生产质量。

\textbf{(4)} 添加约束条件：添加工人每周工作天数、每个班次每天所需人数以及工龄分布比例等约束条件。

\textbf{(5)} 求解模型：使用PuLP库的求解器对模型进行求解，以找到最优的工人排班方案。

通过以上步骤，可以有效地求解出最优的人工排班方案，以最大化生产质量。

\subsubsection{模型输出}

\textbf{排版结果具体见附录一，附录二。}

表 1 数据生成：

\textbf{(1)} 首先，为每个工人创建一个空列表，用于记录其每天的班次或休息状态。

\textbf{(2)} 遍历每一天和每一个工人，通过检查决策变量的值来确定工人在该天是否被安排了班次。

\textbf{(3)} 如果工人被安排了班次，则将班次添加到对应工人的列表中；否则，将"休"添加到列表中表示休息。

\textbf{(4)} 将每个工人的班次情况整理成\(DataFrame\)表格。

\vspace{12pt}
表 2 数据生成：

\textbf{(1)} 创建表格列名，包括日期、班次以及 10 个工人的编号。

\textbf{(2)} 初始化一个空列表用于存储表格数据。

\textbf{(3)} 对于每一天和每一个班次，记录该班次的日期和班次类型。

\textbf{(4)} 根据决策变量的值确定被分配到班次的工人，并确保每个班次有 10 个工人。

\textbf{(5)} 将数据整理成表格形式，包括日期、班次以及 10 个工人的编号。

\vspace{12pt}
数字描述：

\textbf{(1)} 表 1 数据生成的过程可以表示为：

对于每个工人 \( i \)，在每天 \( j \) 中，定义决策变量 \( x_{i,j,k} \) 表示工人 \( i \) 是否被分配到班次 \( k \)。

如果 \( x_{i,j,k} \) 的值为1，则表示工人 \( i \) 在第 \( j \) 天被分配到班次 \( k \)，否则表示工人休息。

表格数据的生成遵循上述规则，并将结果整理成 \(DataFrame\)。

\textbf{(2)} 表 2 数据生成的过程可以表示为：

对于每个班次 \( k \) 和日期 \( j \)，记录每个班次的日期和班次类型。

根据决策变量的值确定被分配到班次的工人，并确保每个班次有 10 个工人。

整理数据成表格形式，包括日期、班次以及 10 个工人的编号。

\subsection{总结}

以上求解过程主要涉及本团队通过机器学习模型精确确定不同工龄段操作人员在整体操作人员中的比例，并运用线性规划模型来优化工人排班方案，从而最大化生产质量。具体而言，本团队研究使用了聚类分析等机器学习模型对操作人员的工龄、工作经验等特征进行聚类，以便更好地理解和描述不同工龄段操作人员的特征和数量分布情况。随后，借助决策树等机器学习模型预测和优化不同工龄段操作人员在不同班次中的分配情况，通过数据驱动的方法提高了排班方案的灵活性和效率。在此基础上，利用线性规划模型构建了一个混合整数线性规划问题，并详细阐述了目标函数、决策变量和约束条件的数学描述。通过求解器对模型进行求解，得出了最优的工人排班方案，以最大化生产质量。本团队采用的研究方法不仅提供了一种有效的工人排班方案优化方法，而且为实际生产中的排班问题提供了数学建模和求解的范例，具有一定的学术和实践意义。

%%%%%%%%%%%%%%%%%%%%%%%%%%%%%%%%%%%%%%%%%%%%%%%%%%%%%%%%%%%%%

\section{模型的评价和推广}

\subsection{模型的优点}
\begin{itemize}[itemindent=2em]
\item 优点1: 利用 \(Matlab\)，\(Python\) ，\(SPSSPRO\) 等软件进行分析，处理问题方便快捷，形象直观。

\item 优点2: \(XGBoost\) 与 \(AdaBoost\) 模型，针对处理含有异常值和稀疏数据的大规模数据集时表现优异，具有较好的鲁棒性，能够显著减少过拟合风险。

\item 优点3: 决策树模型对高维数据的处理能力强，可以直观表示数据之间的关系，易于理解和解释。

\item 优点4: 不受数据分布的影响:\(SVM\) 不依赖数据的概率分布，因此在处理非正态分布数据时依然表现良好。

\item 优点5: 提升整体模型性能,\(Gradient Boosting\)基于接待训练弱分类器处理各类型数据与复杂非线性关系。

\item 优点6: 多算法验证数据特征，采用了\(XGBoost\) 、\(GBDT\) 和 \(RandomForest\) 等算法对各装置故障的数据特征进行验证。这一步骤的创新点在于使用了多种不同的机器学习算法，以验证数据特征的有效性和稳定性。通过比较不同算法的表现，可以更全面地理解数据特征之间的关系，提高了模型的鲁棒性和泛化能力。

\end{itemize}

\subsection{模型的缺点}
\begin{itemize}[itemindent=2em]
\item 缺点1: \(XGBoost\)在处理高度不平衡的数据集时可能表现不佳。
\item 缺点2: \(Gradient Boosting\)模型调参过程相对较复杂，需要较多经验和时间。
\end{itemize}

\subsection{模型的推广}

这些模型可以在多个领域和场景中得到推广和应用：

\textbf{(1)} 决策树模型和\(Gradient Boosting\)

这些模型可以应用于制造业生产线的故障预测和即时报警，帮助企业实现生产线的智能监控和管理。通过分析装置故障数据特征，运用决策树模型和\(Gradient Boosting\)实现自动报警，可以在生产过程中及时发现潜在故障，并采取措施进行修复，以避免生产中断和经济损失。

\textbf{(2)} 聚类分析和\(MILP\)模型

结合聚类分析及\(MILP\)模型，制定出操作人员最佳排班方案，可以帮助企业合理安排人力资源，提高员工工作效率和满意度。通过优化操作人员的排班安排，可以实现生产规`模的扩展而不增加额外的人力成本，从而提高企业的竞争力和盈利能力。

\textbf{(3)} 其他行业的应用

这些模型不仅局限于制造业，在其他行业如电力、交通、医疗等领域也可以得到广泛应用。例如，在电力行业，可以利用这些模型对电网设备的运行状态进行监测和预测，提高电网的可靠性和稳定性；在交通领域，可以利用这些模型对交通流量和拥堵情况进行预测和调控，优化交通运输系统的运行效率和安全性；在医疗领域，可以利用这些模型对患者的病情和治疗效果进行预测和评估，提高医疗服务的质量和效率。

%%%%%%%%%%%%%%%%%%%%%%%%%%%%%%%%%%%%%%%%%%%%%%%%%%%%%%%%%%%%%
%% 参考文献
\newpage

\nocite{*}
\bibliographystyle{gbt7714-numerical}  % 引用格式
\bibliography{ref.bib}  % bib源

%%%%%%%%%%%%%%%%%%%%%%%%%%%%%%%%%%%%%%%%%%%%%%%%%%%%%%%%%%%%%
%% 附录
%\clearpage

%\subsection{Result 4-1}
%\includepdf[pages=-]{result4-1.pdf}

%\subsection{Result 4-2}
%\includepdf[pages=-]{result4-2.pdf}


\end{document}

